\chapter{Perelman's Reduced Volume}
\label{chap:cz-reduced-volume}

\bigskip
In Section 1.5 we introduced the $\mathcal{F}$-functional and the
$\mathcal{W}$-functional of Perelman \cite{P1} and proved their
monotonicity properties under the Ricci flow. In the last section
of the previous chapter we have defined the Li-Yau-Perelman
distance. The main purpose of this chapter is to use the
Li-Yau-Perelman distance to define the Perelman's reduced volume,
which was introduced by Perelman in \cite{P1}, and prove the
monotonicity property of the reduced volume under the Ricci flow.
This new monotonicity formula of Perelman \cite{P1} is more useful
for local considerations, especially when we consider the
formation of singularities in Chapter 6 and work on the Ricci flow
with surgery in Chapter 7. As first applications we will present
two no local collapsing theorems of Perelman \cite{P1} in this
chapter. More applications can be found in Chapter 6 and 7. This
chapter is a detailed exposition of sections 6-8 of Perelman
\cite{P1}.


\section{Riemannian Formalism in Potentially Infinite Dimensions}

In Section 2.6, from an analytic view point, we saw how the Li-Yau
path integral of Perelman's estimate for fundamental solutions to
the conjugate heat equation leads to the Li-Yau-Perelman distance.
In this section we present, from a geometric view point, another
motivation (cf. section 6 of \cite{P1}) how one is lead to the
consideration of the Li-Yau-Perelman distance function, as well as
a reduced volume concept. Interestingly enough, the
Li-Yau-Hamilton quadratic introduced in Section 2.5 appears again
in this geometric consideration.

We consider the Ricci flow
$$
\frac{\partial}{\partial t}g_{ij}=-2R_{ij}
$$
on a manifold $M$ where we assume that $g_{ij}(\cdot,t)$ are
complete and have uniformly bounded curvatures.

Recall from Section 2.5 that the Li-Yau-Hamilton quadratic
introduced in \cite{Ha93} is
$$
Q=M_{ij}W_iW_j+2P_{ijk}U_{ij}W_k+R_{ijkl}U_{ij}U_{kl}
$$
where
\begin{align*}
M_{ij}& = \Delta R_{ij}-\frac{1}{2}\nabla_i\nabla_jR
+2R_{ikjl}R_{kl}-R_{ik}R_{jk}+\frac{1}{2t}R_{ij},\\
P_{ijk}&  = \nabla_iR_{jk}-\nabla_jR_{ik}
\end{align*}
and $U_{ij}$ is any two-form and $W_i$ is any 1-form. Here and
throughout this chapter we do not always bother to raise indices;
repeated indices is short hand for contraction with respect to the
metric.

In \cite{Ha95F}, Hamilton predicted that the Li-Yau-Hamilton
quadratic is some sort of jet extension of positive curvature
operator on some larger space. Such an interpretation of the
Li-Yau-Hamilton quadratic as a curvature operator on the space
$M\times \mathbb{R}^{+}$ was found by Chow and Chu \cite{ChCh95}
where a potentially degenerate Riemannian metric on $M\times
\mathbb{R}^+$ was constructed. The degenerate Riemannian metric on
$M\times \mathbb{R}^+$ is the limit of the following two-parameter
family of Riemannian metrics
$$
g_{N,\delta}(x,t)=g(x,t)+(R(x,t)+\frac{N}{2(t+\delta)})dt^{2}
$$
as $N$ tends to infinity and $\delta$ tends to zero, where
$g(x,t)$ is the solution of the Ricci flow on $M$ and $t\in
\mathbb{R}^{+}$.

To avoid the degeneracy, Perelman \cite{P1} considers the manifold
$\tilde{M}=M\times \mathbb{S}^N\times \mathbb{R}^+$ with the
following metric:
\begin{align*}
\tilde{g}_{ij}&  = g_{ij},\\
\tilde{g}_{\alpha\beta}&  = \tau g_{\alpha\beta},\\
\tilde{g}_{oo}&  = \frac{N}{2\tau}+R,\\
\tilde{g}_{i\alpha} &  = \tilde{g}_{io} =\tilde{g}_{\alpha o}=0,
\end{align*}
where $i,j$ are coordinate indices on $M$; $\alpha,\;\beta$ are
coordinate indices on $\mathbb{S}^N$; and the coordinate $\tau$ on
$R^+$ has index $o$. Let $\tau = T-t$ for some fixed constant $T$.
Then $g_{ij}$ will evolve with $\tau$ by the backward Ricci flow
$\frac{\partial}{\partial \tau}g_{ij}=2R_{ij}$. The metric
$g_{\alpha\beta}$ on $\mathbb{S}^N$ is a metric with constant
sectional curvature $\frac{1}{2N}$.

We remark that the metric $\tilde{g}_{\alpha\beta}$ on
$\mathbb{S}^N$ is chosen so that the product metric
$(\tilde{g}_{ij},\ \tilde{g}_{\alpha\beta})$ on $M\times
\mathbb{S}^N$ evolves by the Ricci flow, while the component
$\tilde{g}_{oo}$ is just the scalar curvature of
$(\tilde{g}_{ij},\ \tilde{g}_{\alpha\beta})$. Thus the metric
$\tilde{g}$ defined on $\tilde{M}=M\times \mathbb{S}^N \times
\mathbb{R}^+$ is exactly a ``regularization" of Chow-Chu's
degenerate metric on $M\times\mathbb{R}^+$.  \vskip 0.2cm

%\noindent{\bf Proposition 3.1.1 } \emph{
% CZ-BASELINE source-line:5122 source-kind:proposition source-id:3.1.1
\begin{proposition}[space-time curvature and the Li–Yau–Hamilton quadratic]
  \label{prop:cz-space-time-curvature-and-the-liyauhamilton-quadratic}
  \dcref{Proposition 3.1.1}
  \group{Cao--Zhu Riemannian Formalism in Potentially Infinite Dimensions}
  \source{cao-zhu}
The components of the curvature tensor of the metric $\tilde{g}$
coincide $($modulo $N^{-1})$ with the components of the
Li-Yau-Hamilton quadratic.
\end{proposition}

%\vskip 0.1cm \noindent{\bf Proof.}
\begin{proof}
By definition, the Christoffel symbols of the metric $\tilde{g}$
are given by the following list: {\allowdisplaybreaks
\begin{align*}
\tilde{\Gamma}_{ij}^k&=\Gamma_{ij}^k,\\
\tilde{\Gamma}_{i\beta}^k&=0\quad \text{and}\quad
\tilde{\Gamma}_{ij}^\gamma=0,\\
\tilde{\Gamma}_{\alpha\beta}^k&=0\quad \text{and}\quad
\tilde{\Gamma}_{i\beta}^\gamma=0,\\
\tilde{\Gamma}_{io}^k&=g^{kl}R_{li}\quad
\text{and}\quad\tilde{\Gamma}_{ij}^o=-\tilde{g}^{oo}R_{ij},\\
\tilde{\Gamma}_{oo}^k&=-\frac{1}{2}g^{kl} \frac{\partial}{\partial
x^l}R\quad
\text{and}\quad\tilde{\Gamma}_{io}^o=\frac{1}{2}\tilde{g}^{oo}
\frac{\partial}{\partial x^i}R,\\
\tilde{\Gamma}_{i\beta}^o&=0,\;\tilde{\Gamma}_{o\beta}^k=0\quad
\text{and}\quad\tilde{\Gamma}_{oj}^\gamma=0,\\
\tilde{\Gamma}_{\alpha\beta}^\gamma&=\Gamma_{\alpha\beta}^\gamma,\\
\tilde{\Gamma}_{\alpha o}^\gamma
&=\frac{1}{2\tau}\delta_\alpha^\gamma\quad
\text{and}\quad\tilde{\Gamma}_{\alpha\beta}^o
=-\frac{1}{2}\tilde{g}^{oo}g_{\alpha\beta},\\
\tilde{\Gamma}_{oo}^\gamma&=0\quad \text{and}\quad
\tilde{\Gamma}_{o\beta}^o=0,\\
\tilde{\Gamma}_{oo}^o&=\frac{1}{2}\tilde{g}^{oo}
\(-\frac{N}{2\tau^2}+\frac{\partial}{\partial\tau}R\).
\end{align*}
} Fix a point $(p,s,\tau)\in M\times \mathbb{S}^N\times
\mathbb{R}^+$ and choose normal coordinates around $p\in M$ and
normal coordinates around $s\in \mathbb{S}^N$ such that
$\Gamma_{ij}^k(p)=0$ and $\Gamma_{\alpha\beta}^\gamma(s)=0$ for
all $i,j,k$ and $\alpha,\beta,\gamma$. We compute the curvature
tensor $\tilde{R}m$ of the metric $\tilde{g}$ at the point as
follows:
\begin{align*}
\tilde{R}_{ijkl}& =R_{ijkl}+\tilde{\Gamma}_{io}^k
\tilde{\Gamma}_{jl}^o-\tilde{\Gamma}_{jo}^k\tilde{\Gamma}_{il}^o=
R_{ijkl}+O\(\frac{1}{N}\),\\
\tilde{R}_{ijk\delta}&  = 0,\\
\tilde{R}_{ij\gamma\delta}&  = 0\quad \text{and} \quad
\tilde{R}_{i\beta
k\delta}=\tilde{\Gamma}_{io}^k\tilde{\Gamma}_{\beta\delta}^o
-\tilde{\Gamma}_{\beta o}^k\tilde{\Gamma}_{i\delta}^o
=-\frac{1}{2}\tilde{g}^{oo}g_{\beta\delta}g^{kl}R_{li}
=O\(\frac{1}{N}\),\\
\tilde{R}_{i\beta\gamma\delta}&  = 0,\\
\tilde{R}_{ijko}&  = \frac{\partial}{\partial
x^i}R_{jk}-\frac{\partial}{\partial
x^j}R_{ik}+\tilde{\Gamma}_{io}^k\tilde{\Gamma}_{jo}^o-
\tilde{\Gamma}_{jo}^k\tilde{\Gamma}_{io}^o=P_{ijk}+O\(\frac{1}{N}\),\\
\tilde{R}_{ioko}&  = -\frac{1}{2}\frac{\partial^2}{\partial
x^i\partial x^k}R-\frac{\partial}{\partial
\tau}(R_{il}g^{lk})+\tilde{\Gamma}_{io}^k\tilde{\Gamma}_{oo}^o
-\tilde{\Gamma}_{oj}^k\tilde{\Gamma}_{io}^j
-\tilde{\Gamma}_{oo}^k\tilde{\Gamma}_{io}^o\\
&  = -\frac{1}{2}\nabla_i\nabla_kR-\frac{\partial}{\partial
\tau}R_{ik}+2R_{ik}R_{lk}-\frac{1}{2\tau}R_{ik}-R_{ij}R_{jk}
+O\(\frac{1}{N}\)\\
&  = M_{ik}+O(\frac{1}{N}),\\
\tilde{R}_{ij\gamma o}&  = 0\quad
\text{and}\quad \tilde{R}_{i\gamma jo}=0,\\
\tilde{R}_{i\beta\gamma o}&  = -\tau\tilde{\Gamma}_{\beta
o}^\gamma\tilde{\Gamma}_{io}^o=O\(\frac{1}{N}\)\quad
\mbox{and}\quad
\tilde{R}_{io\gamma\delta}=0,\\
\tilde{R}_{io\gamma o}&  = 0,\\
\tilde{R}_{\alpha\beta\gamma o}&  = 0, \\
\tilde{R}_{\alpha o\gamma o}& =
\(\frac{1}{2\tau^2}\delta_\alpha^\gamma +\tilde{\Gamma}_{\alpha
o}^\gamma\tilde{\Gamma}_{oo}^o
-\tilde{\Gamma}_{o\beta}^\gamma\tilde{\Gamma}_{\alpha
o}^\beta\)\tau=O\(\frac{1}{N}\),\\
\tilde{R}_{\alpha\beta\gamma\delta}&  = O\(\frac{1}{N}\).
\end{align*}

Thus the components of the curvature tensor of the metric
$\tilde{g}$ coincide (modulo $N^{-1}$) with the components of the
Li-Yau-Hamilton quadratic.
%$$\eqno \#$$ \vskip 0.3cm
\end{proof}

The following observation due to Perelman \cite{P1} gives an
important motivation to define Perelman's reduced volume.

%0.2cm \noindent{\bf Corollary 3.1.2}\ \emph{
% CZ-BASELINE source-line:5214 source-kind:corollary source-id:3.1.2
\begin{corollary}[space-time Ricci flatness modulo inverse dimension]
  \label{cor:cz-space-time-ricci-flatness-modulo-inverse-dimension}
  \dcref{Corollary 3.1.2}
  \group{Cao--Zhu Riemannian Formalism in Potentially Infinite Dimensions}
  \source{cao-zhu}
All components of the Ricci tensor of $\tilde{g}$ are zero
$($modulo $N^{-1})$.
\end{corollary}

%\vskip 0.1cm\noindent{\bf Proof.}\
\begin{proof}
{}From the list of the components of the curvature tensor of
$\tilde{g}$ given above, we have
\begin{align*}
\tilde{R}_{ij}& =
\tilde{g}^{kl}\tilde{R}_{ijkl}+\tilde{g}^{\alpha\beta}\tilde{R}_{i\alpha
j\beta}+\tilde{g}^{oo}\tilde{R}_{iojo}\\
& =R_{ij}-\frac{1}{2\tau}g^{\alpha\beta}\tilde{g}^{oo}
g^{\alpha\beta}R_{ij}+\tilde{g}^{oo}
\(M_{ij}-\frac{1}{2\tau}R_{ij}+O\(\frac{1}{N}\)\)\\
&  =R_{ij}-\frac{N}{2\tau}\tilde{g}^{oo}R_{ij}+O\(\frac{1}{N}\)\\
&  =O\(\frac{1}{N}\),\\
\tilde{R}_{i\gamma}&  = \tilde{g}^{kl}\tilde{R}_{ik\gamma
l}+\tilde{g}^{\alpha\beta}\tilde{R}_{i\alpha\gamma\beta}
+\tilde{g}^{oo}\tilde{R}_{io\gamma o}=0,\\
\tilde{R}_{io}& =\tilde{g}^{kl}\tilde{R}_{ikol}
+\tilde{g}^{\alpha\beta}\tilde{R}_{i\alpha
o\beta}+\tilde{g}^{oo}\tilde{R}_{iooo}\\
&  = -g^{kl}P_{ikl}+O\(\frac{1}{N}\),\\
\tilde{R}_{\alpha\beta}&  = \tilde{g}^{kl}\tilde{R}_{\alpha k\beta
l}+\tilde{g}^{\gamma\delta}\tilde{R}_{\alpha\gamma\beta\delta}
+\tilde{g}^{oo}\tilde{R}_{\alpha o \beta o}\\
&  = O\(\frac{1}{N}\),\\
\tilde{R}_{\alpha o}&  = \tilde{g}^{kl}\tilde{R}_{\alpha
kol}+\tilde{g}^{\beta\gamma}\tilde{R}_{\alpha\beta
o\gamma}+\tilde{g}^{oo}\tilde{R}_{\alpha ooo}=0,\\
\tilde{R}_{oo}& =
\tilde{g}^{kl}\tilde{R}_{okol}+\tilde{g}^{\alpha\beta}\tilde{R}_{o\alpha
o\beta}+\tilde{g}^{oo}\tilde{R}_{oooo}\\
&  = g^{kl}\(M_{kl}+O\(\frac{1}{N}\)\)+O\(\frac{1}{N}\).
\end{align*}
Since $\tilde{g}^{oo}$ is of order $N^{-1}$, we see that the norm
of the Ricci tensor is given by
$$
|\widetilde{\Ric}|_{\tilde{g}}=O\(\frac{1}{N}\).
$$
This proves the result.
\end{proof}
%$$\eqno \#$$ \vskip 0.3cm

We now use the Ricci-flatness of the metric $\tilde{g}$ to
interpret the Bishop-Gromov relative volume comparison theorem
which will motivate another monotonicity formula for the Ricci
flow. The argument in the following will not be rigorous. However
it gives an intuitive picture of what one may expect. Consider a
metric ball in $(\tilde{M},\tilde{g})$ centered at some point
$(p,s,0)\in\tilde{M}$. Note that the metric of the sphere
$\mathbb{S}^N$ at $\tau=0$ degenerates and it shrinks to a point.
Then the shortest geodesic $\gamma(\tau)$ between $(p,s,0)$ and an
arbitrary point $(q,\bar{s},\bar\tau)\in\tilde{M}$ is always
orthogonal to the $\mathbb{S}^N$ fibre. The length of
$\gamma(\tau)$ can be computed as
\begin{align*}
& \int_0^{\bar\tau}\sqrt{\(\frac{N}{2\tau}+R\)
+|\dot{\gamma}(\tau)|^2_{g_{ij}(\tau)}}d\tau\\
&
=\sqrt{2N\bar\tau}+\frac{1}{\sqrt{2N}}\int_0^{\bar\tau}\sqrt{\tau}
(R+|\dot{\gamma}(\tau)|^2_{g_{ij}})d\tau+O(N^{-\frac{3}{2}}).
\end{align*}
Thus a shortest geodesic should minimize
$$
\mathcal{L}(\gamma)=\int_0^{\bar\tau}\sqrt{\tau}
(R+|\dot{\gamma}(\tau)|^2_{g_{ij}})d\tau.
$$

Let $L(q,\bar\tau)$ denote the corresponding minimum. We claim
that a metric sphere $S_{\tilde{M}}(\sqrt{2N\bar\tau})$ in
$\tilde{M}$ of radius $\sqrt{2N\bar\tau}$ centered at $(p,s,0)$ is
$O(N^{-1})$-close to the hypersurface $\{\tau=\bar\tau \}$.
Indeed, if $(x,s',\tau(x))$ lies on the metric sphere
$S_{\tilde{M}}(\sqrt{2N\bar\tau})$, then the distance between
$(x,s',\tau(x))$ and $(p,s,0)$ is
$$
\sqrt{2N\bar\tau}
=\sqrt{2N\tau(x)}+\frac{1}{\sqrt{2N}}L(x,\tau(x))+O\(N^{-\frac{3}{2}}\)
$$
which can be written as
$$
\sqrt{\tau(x)}-\sqrt{\bar\tau}=-\frac{1}{2N}L(x,\tau(x))+O(N^{-2})
=O(N^{-1}).$$ This shows that the metric sphere
$S_{\tilde{M}}(\sqrt{2N\bar\tau})$ is $O(N^{-1})$-close to the
hypersurface $\{\tau=\bar\tau \}$. Note that the metric
$g_{\alpha\beta}$ on $\mathbb{S}^N$ has constant sectional
curvature $\frac{1}{2N}$. Thus
\begin{align*}
& {\rm Vol}\,\(S_{\tilde{M}}\(\sqrt{2N\bar\tau}\)\) \\
& \approx \int_{M}\(\int_{S^N}dV_{\tau(x)g_{\alpha\beta}}\)dV_{g_{ij}}(x)\\
& = \int_M(\tau(x))^{\frac{N}{2}}{\rm Vol}\,(S^N)dV_M\\
& \approx (2N)^{\frac{N}{2}}\omega_N\int_M\(\sqrt{\bar\tau}
-\frac{1}{2N}L(x,\tau(x))+O(N^{-2})\)^NdV_M\\
& \approx (2N)^{\frac{N}{2}}\omega_N\int_M
\(\sqrt{\bar\tau}-\frac{1}{2N}L(x,\bar\tau)+o(N^{-1})\)^NdV_M,
\end{align*}
where $\omega_N$ is the volume of the standard $N$-dimensional
sphere. Now the volume of Euclidean sphere of radius
$\sqrt{2N\bar\tau}$ in $\mathbb{R}^{n+N+1}$ is
$$
{\rm Vol}\,(S_{\mathbb{R}^{n+N+1}}(\sqrt{2N\bar\tau}))
=(2N\bar\tau)^\frac{N+n}{2}\omega_{n+N}.
$$
Thus we have
$$
\frac{{\rm Vol}\,(S_{\tilde{M}}(\sqrt{2N\bar\tau}))}{{\rm
Vol}\,(S_{
\mathbb{R}^{n+N+1}}(\sqrt{2N\bar\tau}))}\approx\text{const}\cdot
N^{-\frac{n}{2}}\cdot \int_M(\bar\tau)^{-\frac{n}{2}}\exp
\left\{{-\frac{1}{2\sqrt{\bar\tau}}L(x,\bar\tau)}\right\}dV_M.
$$
Since the Ricci curvature of $\tilde{M}$ is zero (modulo
$N^{-1}$), the Bishop-Gromov volume comparison theorem then
suggests that the integral
$$
\tilde{V}(\bar\tau)\stackrel{\Delta}{=}\int_M
(4\pi\bar\tau)^{-\frac{n}{2}}\exp
\left\{-\frac{1}{2\sqrt{\bar\tau}}L(x,\bar\tau)\right\}dV_M,
$$
which we will call {\bf Perelman's reduced
volume}\index{Perelman's reduced volume}, should be nonincreasing
in $\bar\tau$. A rigorous proof of this monotonicity property will
be given in the next section. One should note the analog of
reduced volume with the heat kernel and there is a parallel
calculation for the heat kernel of the Shr\"odinger equation in
the paper of Li-Yau \cite{LY}.


\section{Comparison Theorems for Perelman's Reduced Volume}

In this section we will write the Ricci flow in the backward
version
$$
\frac{\partial}{\partial \tau}g_{ij}=2R_{ij}
$$
on a manifold $M$ with $\tau=\tau(t)$ satisfying $d\tau/dt=-1$ (in
practice we often take $\tau=t_0-t$ for some fixed time $t_0$). We
always assume that either $M$ is compact or $g_{ij}(\tau)$ are
complete and have uniformly bounded curvature. To each (smooth)
space curve $\gamma(\tau)$, $0<\tau_1\le\tau\le\tau_2$, in $M$, we
define its {\bf $\mathcal{L}$-length}\index{$\mathcal{L}$-length}
as
$$
\mathcal{L}(\gamma)
=\int_{\tau_1}^{\tau_2}\sqrt{\tau}(R(\gamma(\tau),\tau)+
|\dot{\gamma}(\tau)|^2_{g_{ij}(\tau)})d\tau.
$$

Let $X(\tau)=\dot{\gamma}(\tau)$, and let $Y(\tau)$ be any
(smooth) vector field along $\gamma(\tau)$. First of all, we
compute the first variation formula for $\mathcal{L}$-length (cf.
section 7 of \cite{P1}).

%\vskip 0.3cm\noindent{\bf Lemma 3.2.1}
% CZ-BASELINE source-line:5371 source-kind:lemma source-id:3.2.1
\begin{lemma}[first variation of L-length]
  \label{lem:cz-first-variation-of-l-length}
  \dcref{Lemma 3.2.1}
  \group{Cao--Zhu Comparison Theorems for Perelman's Reduced Volume}
  \source{cao-zhu}
$$
\delta_Y(\mathcal{L})=2\sqrt{\tau}\<X,Y\>|_{\tau_1}^{\tau_2}
+\int_{\tau_1}^{\tau_2}\sqrt{\tau}\left\<Y,\nabla
R-2\nabla_XX-4\Ric(\cdot,X)-\frac{1}{\tau}X\right\>d\tau
$$
where $\<\cdot,\cdot\>$ denotes the inner product with respect to
the metric $g_{ij}(\tau)$.
\end{lemma}

%\vskip 0.1cm\noindent{\bf Proof.}\
\begin{proof}
By direct computations,
\begin{align*}
\delta_Y(\mathcal{L}) &
=\int_{\tau_1}^{\tau_2}\sqrt{\tau}(\<\nabla R,Y\>
+2\<X,\nabla_YX\>)d\tau\\
&  = \int_{\tau_1}^{\tau_2}\sqrt{\tau}(\<\nabla R,Y\>
+2\<X,\nabla_XY\>)d\tau\\
&  = \int_{\tau_1}^{\tau_2}\sqrt{\tau}\(\<\nabla R,Y\>
+2\frac{d}{d\tau}\<X,Y\>-2\<\nabla_XX,Y\>-4\Ric(X,Y)\)d\tau\\
& =2\sqrt{\tau}\<X,Y\>|_{\tau_1}^{\tau_2}
+\int_{\tau_1}^{\tau_2}\sqrt{\tau}\left\<Y,\nabla R
-2\nabla_XX-4\Ric(\cdot,X)-\frac{1}{\tau}X\right\>d\tau.
\end{align*}
%$$\eqno \#$$ \vskip 0.3cm
\end{proof}

A smooth curve $\gamma(\tau)$ in $M$ is called an {\bf
$\mathcal{L}$-geodesic}\index{$\mathcal{L}$-geodesic}\index{$\mathcal{L}$-geodesic!curve}
if it satisfies the following {\bf $\mathcal{L}$-geodesic
equation}\index{$\mathcal{L}$-geodesic!equation} \be
\nabla_XX-\frac{1}{2}\nabla R+\frac{1}{2\tau}X+2\Ric(X,\cdot)=0.
%\eqno(3.2.1)
\ee Given any two points $p,q\in M$ and $\tau_2>\tau_1>0$, there
always exists an $\mathcal{L}$-shortest curve (or shortest
$\mathcal{L}$-geodesic) $\gamma(\tau)$:
$[\tau_1,\tau_2]\rightarrow M$ connecting $p$ to $q$ which
satisfies the above $\mathcal{L}$-geodesic equation. Multiplying
the $\mathcal{L}$-geodesic equation (3.2.1) by $\sqrt{\tau}$, we
get
$$
\nabla_X(\sqrt{\tau}X)=\frac{\sqrt{\tau}}{2}\nabla R
-2\sqrt{\tau}\Ric(X,\cdot)\quad \text{on}\quad [\tau_1,\tau_2],
$$
or equivalently
$$
\frac{d}{d\tau}(\sqrt{\tau}X)=\frac{\sqrt{\tau}}{2}\nabla R
-2\Ric(\sqrt{\tau}X,\cdot) \quad \text{on}\quad [\tau_1,\tau_2].
$$
Thus if a continuous curve, defined on $[0,\tau_2]$, satisfies the
$\mathcal{L}$-geodesic equation on every subinterval
$0<\tau_1\le\tau\le\tau_2$, then $\sqrt{\tau_1}X(\tau_1)$ has a
limit as $\tau_1\rightarrow 0^+$. This allows us to extend the
definition of the $\mathcal{L}$-length to include the case
$\tau_1=0$ for all those (continuous) curves $\gamma:
[0,\tau_2]\rightarrow M$ which are smooth on $(0,\tau_2]$ and have
limits
$\lim\limits_{\tau\rightarrow0^+}\sqrt{\tau}\dot{\gamma}(\tau)$.
Clearly, there still exists an $\mathcal{L}$-shortest curve
$\gamma(\tau):[0,\tau_2]\rightarrow M$ connecting arbitrary two
points $p,q\in M$ and satisfying the $\mathcal{L}$-geodesic
equation (3.2.1) on $(0,\tau_2]$. Moreover, for any vector $v \in
T_pM$, we can find an $\mathcal{L}$-geodesic $\gamma(\tau)$
starting at $p$ with
$\lim\limits_{\tau\rightarrow0^+}\sqrt{\tau}\dot{\gamma}(\tau)=v$.

{}From now on, we fix a point $p\in M$ and set $\tau_1=0$. The
{\bf $\mathcal{L}$-distance function}\index{$\mathcal{L}$-distance
function} on the space-time $M\times \mathbb{R}^+$ is denoted by
$L(q,\bar{\tau})$ and defined to be the $\mathcal{L}$-length of
the $\mathcal{L}$-shortest curve $\gamma(\tau)$ connecting $p$ and
$q$ with $0\le\tau\le\bar{\tau}$.

Consider a shortest $\mathcal{L}$-geodesic $\gamma: [0,\bar \tau]
\rightarrow M$ connecting $p$ to $q$. In the computations below we
pretend that $\mathcal{L}$-shortest geodesics between $p$ and $q$
are unique for all pairs $(q,\bar \tau)$; if this is not the case,
the inequalities that we obtain are still valid, by a standard
barrier argument, when understood in the sense of distributions
(see, for example, \cite{ScY}).

The first variation formula in Lemma 3.2.1 implies that
$$
\nabla_YL(q,\bar{\tau}) =\left\<2\sqrt{\bar{\tau}}X(\bar{\tau}),
Y(\bar{\tau})\right\>.
$$
Thus
$$
\nabla L(q,\bar{\tau}) =2\sqrt{\bar{\tau}}X(\bar{\tau}),
$$
and \be |\nabla L|^2 =4\bar{\tau}|X|^2
=-4\bar{\tau}R+4\bar{\tau}(R+|X|^2).%\eqno(3.2.2)
\ee We also compute
\begin{align}
L_{\bar{\tau}}(\gamma(\bar{\tau}),\bar{\tau})& =
\frac{d}{d\tau}L(\gamma(\tau),\tau)|_{\tau=\bar{\tau}}
-\<\nabla L,X\>\\
&  = \sqrt{\bar{\tau}}(R+|X|^2)-2\sqrt{\bar{\tau}}|X|^2 \nn\\
&  = 2\sqrt{\bar{\tau}}R-\sqrt{\bar{\tau}}(R+|X|^2).\nn
%(3.2.3)
\end{align}
To evaluate $R+|X|^2$, we compute by using (3.2.1),
\begin{align*}
& \frac{d}{d\tau}(R(\gamma(\tau),\tau)+|X(\tau)|^2_{g_{ij}(\tau)})\\
& = R_\tau+\<\nabla R,X\>+2\<\nabla_XX,X\>+2{ \Ric}(X,X)\\
& = R_\tau+\frac{1}{\tau}R+2\<\nabla R,X\>-2 {\Ric}(X,X)
-\frac{1}{\tau}(R+|X|^2)\\
&  = -Q(X)-\frac{1}{\tau}(R+|X|^2),
\end{align*}
where
$$
Q(X)=-R_\tau-\frac{R}{\tau}-2\<\nabla R,X\>+2{ \Ric}(X,X)
$$
is the trace Li-Yau-Hamilton quadratic in Corollary 2.5.5. Hence
\begin{align*}
\frac{d}{d\tau}(\tau^\frac{3}{2}(R+|X|^2))|_{\tau=\bar{\tau}} &
=\frac{1}{2}\sqrt{\bar{\tau}}(R+|X|^2)
-\bar{\tau}^\frac{3}{2}Q(X)\\
&
=\frac{1}{2}\frac{d}{d\tau}L(\gamma(\tau),\tau)|_{\tau=\bar{\tau}}
-\bar{\tau}^\frac{3}{2}Q(X).
\end{align*}
Therefore, \be
\bar{\tau}^\frac{3}{2}(R+|X|^2)=\frac{1}{2}L(q,\bar{\tau})-K,
%\eqno(3.2.4)
\ee where \be K=\int_0^{\bar{\tau}}\tau^\frac{3}{2}Q(X)d\tau.
%\eqno(3.2.5)
\ee Combining (3.2.2) with (3.2.3), we obtain \be |\nabla L|^2
=-4\bar{\tau}R+\frac{2}{\sqrt{\bar{\tau}}}L
-\frac{4}{\sqrt{\bar{\tau}}}K  %\eqno(3.2.6)
\ee and \be L_{\bar{\tau}}
=2\sqrt{\bar{\tau}}R-\frac{1}{2\bar{\tau}}L
+\frac{1}{\bar{\tau}}K.  % \eqno(3.2.7)
\ee

Next we compute the second variation of an $\mathcal{L}$-geodesic
(cf. section 7 of \cite{P1}).

%\noindent{\bf Lemma 3.2.2}
% CZ-BASELINE source-line:5511 source-kind:lemma source-id:3.2.2
\begin{lemma}[second variation of L-length]
  \label{lem:cz-second-variation-of-l-length}
  \dcref{Lemma 3.2.2}
  \group{Cao--Zhu Comparison Theorems for Perelman's Reduced Volume}
  \source{cao-zhu}
For any $\mathcal{L}$-geodesic $\gamma$, we have
\begin{align*}
\delta_Y^2(\mathcal{L}) &
=2\sqrt{\tau}\<\nabla_YY,X\>|_0^{\bar{\tau}}
+\int_0^{\bar{\tau}}\sqrt{\tau}[2|\nabla_XY|^2
+2\<R(Y,X)Y,X\>\\
&\quad +\nabla_Y\nabla_YR+2\nabla_X{\Ric}(Y,Y)-4\nabla_Y
\Ric(Y,X)]d\tau.
\end{align*}
\end{lemma}

%\vskip 0.1cm\noindent{\bf Proof.}\
\begin{proof}
We compute
\begin{align*}
\delta_Y^2(\mathcal{L})
& = Y\(\int_0^{\bar{\tau}}\sqrt{{\tau}}(Y(R)+2\<\nabla_YX,X\>)d\tau\)\\
& = \int_0^{\bar{\tau}}\sqrt{{\tau}}(Y(Y(R))
+2\<\nabla_Y\nabla_YX,X\>+2|\nabla_YX|^2)d\tau\\
& =\int_0^{\bar{\tau}}\sqrt{{\tau}}(Y(Y(R))
+2\<\nabla_Y\nabla_XY,X\>+2|\nabla_XY|^2)d\tau
\end{align*}
and
\begin{align*}
& 2\<\nabla_Y\nabla_XY,X\>\\
&  = 2\<\nabla_X\nabla_YY,X\>+2\<R(Y,X)Y,X\>\\
&  = 2\frac{d}{d\tau}\<\nabla_YY,X\>-4{ \Ric}(\nabla_YY,X)
-2\<\nabla_YY,\nabla_XX\> \\
&\quad-\(2\left\<\frac{d}{d\tau}\nabla_YY,X\right\>
-2\<\nabla_X\nabla_YY,X\>\)+2\<R(Y,X)Y,X\>\\
&  = 2\frac{d}{d\tau}\<\nabla_YY,X\>
-4{\Ric}(\nabla_YY,X)-2\<\nabla_YY,\nabla_XX\>\\
&\quad -2\left\<Y^iY^j(g^{kl}(\nabla_iR_{lj}+\nabla_jR_{li}
-\nabla_lR_{ij}))\frac{\partial}{\partial x^k},X\right\>+2\<R(Y,X)Y,X\>\\
& = 2\frac{d}{d\tau}\<\nabla_YY,X\>-4{\Ric}(\nabla_YY,X)
-2\<\nabla_YY,\nabla_XX\>-4\nabla_Y{\Ric}(X,Y)\\
&\quad +2\nabla_X{\Ric}(Y,Y)+2\<R(Y,X)Y,X\>,
\end{align*}
where we have used the computation $$\frac{\partial}{\partial
\tau}\Gamma_{ij}^k=g^{kl}(\nabla_iR_{lj}+\nabla_jR_{li}-\nabla_lR_{ij}).$$
Thus by using the $\mathcal{L}$-geodesic equation (3.2.1), we get
\begin{align*}
\delta_Y^2(\mathcal{L}) &
=\int_0^{\bar{\tau}}\sqrt{\tau}\bigg[Y(Y(R))
+2\frac{d}{d\tau}\<\nabla_YY,X\>-4\Ric(\nabla_YY,X)\\
&\quad-2\<\nabla_YY,\nabla_XX\>
-4\nabla_Y\Ric(X,Y)+2\nabla_X \Ric(Y,Y) \\
&\quad+2\<R(Y,X)Y,X\>+2|\nabla_XY|^2\bigg]d\tau\\
&  =
\int_0^{\bar{\tau}}\bigg[2\sqrt{\tau}\frac{d}{d\tau}\<\nabla_YY,X\>
+\frac{1}{\sqrt{\tau}}\<\nabla_YY,X\>\bigg]d\tau\\
&\quad +\int_0^{\bar{\tau}}\sqrt{\tau}[Y(Y(R))
-\<\nabla_YY,\nabla R\>-4\nabla_Y \Ric(X,Y)\\
&\quad +2\nabla_X \Ric(Y,Y)+2\<R(Y,X)Y,X\>+2|\nabla_XY|^2]d\tau\\
&  = 2\sqrt{\tau}\<\nabla_YY,X\>|_0^{\bar{\tau}}
+\int_0^{\bar{\tau}}\sqrt{\tau}[2|\nabla_XY|^2
+2\<R(Y,X)Y,X\> \\
&\quad+\nabla_Y\nabla_YR -4\nabla_Y \Ric(X,Y) +2\nabla_X
\Ric(Y,Y)]d\tau.
\end{align*}

\vspace*{-5mm}
\end{proof}
%$$\eqno\#$$\vskip 0.3cm

We now use the above second variation formula to estimate the
Hessian of the $\mathcal{L}$-distance function.

Let $\gamma(\tau):[0,\bar{\tau}]\rightarrow M$ be an
$\mathcal{L}$-shortest curve connecting $p$ and $q$ so that the
$\mathcal{L}$-distance function $L=L(q,\bar\tau)$ is given by the
$\mathcal{L}$-length of $\gamma$. We fix a vector $Y$ at
$\tau=\bar{\tau}$ with $|Y|_{g_{ij}(\bar \tau)}=1$, and extend $Y$
along the $\mathcal{L}$-shortest geodesic $\gamma$ on
$[0,\bar{\tau}]$ by solving the following ODE \be
\nabla_XY=-\Ric(Y,\cdot)+\frac{1}{2\tau}Y.  %\eqno(3.2.8)
\ee This is similar to the usual parallel translation and
multiplication with proportional parameter. Indeed, suppose
$\{Y_1,\ldots,Y_n\}$ is an orthonormal basis at $\tau=\bar{\tau}$
(with respect to the metric $g_{ij}(\bar \tau)$) and extend this
basis along the $\mathcal{L}$-shortest geodesic $\gamma$ by
solving the above ODE (3.2.8). Then
\begin{align*}
\frac{d}{d\tau}\<Y_i,Y_j\>
&  = 2\Ric(Y_i,Y_j)+\<\nabla_XY_i,Y_j\>+\<Y_i,\nabla_XY_j\>\\
&  = \frac{1}{\tau}\<Y_i,Y_j\>
\end{align*}
for all $i,j$. Hence, \be \<Y_i(\tau),Y_j(\tau)\>
=\frac{\tau}{\bar{\tau}}\delta_{ij}%\eqno(3.2.9)
\ee and $\{Y_1(\tau),\ldots,Y_n(\tau)\}$ remains orthogonal on
$[0,\bar{\tau}]$ with $Y_i(0)=0,\;i=1,\ldots,n$.

%\vskip 0.2cm \noindent{\bf Proposition 3.2.3}\ \emph{
% CZ-BASELINE source-line:5605 source-kind:proposition source-id:3.2.3
\begin{proposition}[L-Jacobi comparison inequality]
  \label{prop:cz-l-jacobi-comparison-inequality}
  \dcref{Proposition 3.2.3}
  \group{Cao--Zhu Comparison Theorems for Perelman's Reduced Volume}
  \source{cao-zhu}
Given any unit vector $Y$ at any point $q \in M$ with
$\tau=\bar{\tau}$, consider an $\mathcal{L}$-shortest geodesic
$\gamma$ connecting $p$ to $q$ and extend $Y$ along $\gamma$ by
solving the {\rm ODE} $(3.2.8)$. Then the Hessian of the
$\mathcal{L}$-distance function $L$ on $M$ with $\tau=\bar{\tau}$
satisfies
$$
{\rm
Hess}_L(Y,Y)\leq\frac{1}{\sqrt{\bar{\tau}}}-2\sqrt{\bar{\tau}}\Ric(Y,Y)
-\int_0^{\bar{\tau}}\sqrt{{\tau}}Q(X,Y)d\tau
$$
in the sense of distributions, where
\begin{align*}
Q(X,Y)&  = -\nabla_Y\nabla_YR-2\<R(Y,X)Y,X\>
-4\nabla_X \Ric(Y,Y)+4\nabla_Y \Ric(Y,X)\\
&\quad -2\Ric_\tau(Y,Y)+2|\Ric(Y,\cdot)|^2-\frac{1}{\tau}\Ric(Y,Y)
\end{align*}
is the Li-Yau-Hamilton quadratic. Moreover the equality holds if
and only if the vector field $Y(\tau),\tau\in[0,\bar{\tau}]$, is
an {\bf $\mathcal{L}$-Jacobian field} $($i.e., $Y$ is the
derivative of a variation of $\gamma$ by
$\mathcal{L}$-geodesics$)$.\index{$\mathcal{L}$-Jacobian field}
\end{proposition}

%\vskip 0.1cm\noindent{\bf Proof.}
\begin{proof}
  \uses{lem:cz-second-variation-of-l-length}
As said before, we pretend that the shortest
$\mathcal{L}$-geodesics between $p$ and $q$ are unique so that
$L(q,\bar{\tau})$ is smooth. Otherwise, the inequality is still
valid, by a standard barrier argument, when understood in the
sense of distributions (see, for example, \cite{ScY}).

Recall that $\nabla L(q,\bar{\tau})=2\sqrt{\bar{\tau}}X$. Then
$\<\nabla_YY,\nabla L\>=2\sqrt{\bar{\tau}}\<\nabla_YY,X\>.$ We
compute by using Lemma 3.2.2, (3.2.8) and (3.2.9),
\begin{align*}
{\rm Hess}_L(Y,Y)&  = Y(Y(L))(\bar{\tau})-\langle\nabla_YY,\nabla
L\rangle(\bar\tau)\\
& \leq \delta_Y^2({\mathcal{L}})-2\sqrt{\bar{\tau}}
\langle\nabla_YY,X\rangle(\bar\tau)\\
&  = \int_0^{\bar{\tau}}\sqrt{\tau}[2 |\nabla_XY|^2
+2\langle R(Y,X)Y,X\rangle+\nabla_Y\nabla_YR\\
&\quad +2\nabla_X{\Ric}(Y,Y)-4\nabla_Y{\Ric}(Y,X)]d\tau\\
&  = \int_0^{\bar{\tau}}\!\!\sqrt{\tau}\bigg[2
\bigg|-\Ric(Y,\cdot)
+\frac{1}{2\tau}Y\bigg|^2\!+2\langle R(Y,X)Y,X\rangle+\nabla_Y\nabla_YR\\
&\quad +2\nabla_X \Ric(Y,Y)-4\nabla_Y \Ric(Y,X)\bigg]d\tau\\
& =\int_0^{\bar{\tau}}\sqrt{\tau}\bigg[2|\Ric(Y,\cdot)|^2
-\frac{2}{\tau}\Ric(Y,Y)
+\frac{1}{2\tau\bar{\tau}}+2\langle R(Y,X)Y,X\rangle\\
&\quad +\nabla_Y\nabla_YR+2\nabla_X \Ric(Y,Y)-4\nabla_Y
\Ric(Y,X)\bigg]d\tau.
\end{align*}
Since
\begin{align*}
\frac{d}{d\tau}\Ric(Y,Y)
&  =\Ric_\tau(Y,Y)+\nabla_X \Ric(Y,Y)+2\Ric(\nabla_XY,Y)\\
&  = \Ric_\tau(Y,Y)+\nabla_X
\Ric(Y,Y)-2|\Ric(Y,\cdot)|^2+\frac{1}{\tau}\Ric(Y,Y),
\end{align*}
we have
\begin{align*}
&{\rm Hess}_L(Y,Y) \\
& \leq\int_0^{\bar{\tau}}\sqrt{\tau}\bigg[2|\Ric(Y,\cdot)|^2
-\frac{2}{\tau}\Ric(Y,Y)+\frac{1}{2\tau\bar{\tau}}
+2\langle R(Y,X)Y,X\rangle\\
&\quad +\nabla_Y\nabla_YR-4(\nabla_Y \Ric)(X,Y)
-\bigg(2\frac{d}{d\tau}\Ric(Y,Y)-2\Ric_\tau(Y,Y)\\
&\quad +4|\Ric(Y,\cdot)|^2-\frac{2}{\tau}\Ric(Y,Y)\bigg)
+4\nabla_X \Ric(Y,Y)\bigg]d\tau\\
& =-\int_0^{\bar{\tau}}\bigg[2\sqrt{\tau}\frac{d}{d\tau}\Ric(Y,Y)
+\frac{1}{\sqrt{\tau}}\Ric(Y,Y)\bigg]d\tau+\frac{1}{2\bar\tau}
\int_0^{\bar{\tau}}\frac{1}{\sqrt{\tau}}d\tau\\
&\quad +\int_0^{\bar{\tau}}\sqrt{\tau}\bigg[2\langle
R(Y,X)Y,X\rangle
+\nabla_Y\nabla_YR+\frac{1}{\tau}\Ric(Y,Y)\\
&\quad +4(\nabla_X \Ric(Y,Y)-\nabla_Y \Ric(X,Y))+2\Ric_\tau(Y,Y)
-2|\Ric(Y,\cdot)|^2\bigg]d\tau\\
& =\frac{1}{\sqrt{\bar{\tau}}}-2\sqrt{\bar{\tau}}\Ric(Y,Y)
-\int_0^{\bar{\tau}}\sqrt{\tau}Q(X,Y)d\tau.
\end{align*}
This proves the inequality.

As usual, the quadratic form
$$
I(V,V)=\int_0^{\bar{\tau}}\sqrt{\tau}[2 |\nabla_XV|^2 +2\langle
R(V,X)V,X\rangle+\nabla_V\nabla_VR
$$
$$
+2\nabla_X \Ric(V,V)-4\nabla_V \Ric(V,X)]d\tau,
$$
for any vector field $V$ along $\gamma$, is called the index form.
Since $\gamma$ is shortest, the standard Dirichlet principle for
$I(V,V)$ implies that the equality holds if and only if the vector
field $Y$ is the derivative of a variation of $\gamma$ by
$\mathcal{L}$-geodesics.
\end{proof}
%$$\eqno\#$$

%\vskip 0.3cm \noindent {\bf Corollary 3.2.4}\  \emph{
% CZ-BASELINE source-line:5706 source-kind:corollary source-id:3.2.4
\begin{corollary}[distributional comparison for reduced distance]
  \label{cor:cz-distributional-comparison-for-reduced-distance}
  \dcref{Corollary 3.2.4}
  \group{Cao--Zhu Comparison Theorems for Perelman's Reduced Volume}
  \source{cao-zhu}
We have
$$
\Delta L\leq\frac{n}{\sqrt{\bar{\tau}}}
-2\sqrt{\bar{\tau}}R-\frac{1}{\bar{\tau}}K
$$
in the sense of distribution. Moreover, the equality holds if and
only if we are on a gradient shrinking soliton with
$$
R_{ij}+\frac{1}{2\sqrt{\bar\tau}}\nabla_i\nabla_jL
=\frac{1}{2\bar\tau}g_{ij}.$$
\end{corollary}

%\vskip 0.1cm \noindent {\bf Proof.}\ \ \
\begin{proof}
  \uses{prop:cz-l-jacobi-comparison-inequality}
Choose an orthonormal basis $\{Y_1,\ldots,Y_n\}$ at
$\tau=\bar{\tau}$ and extend them along the shortest
$\mathcal{L}$-geodesic $\gamma$ to get vector fields $Y_i(\tau)$,
$i=1,\ldots,n$, by solving the ODE (3.2.8), with $\langle
Y_i(\tau),Y_j(\tau)\rangle=\frac{\tau}{\bar{\tau}}\delta_{ij}$ on
$[0,\bar{\tau}]$. Taking $Y=Y_i$ in Proposition 3.2.3 and summing
over $i$, we get
\begin{align}
\Delta L& \leq
\frac{n}{\sqrt{\bar{\tau}}}-2\sqrt{\bar{\tau}}R-\sum\limits_{i=1}^n
\int_0^{\bar{\tau}}\sqrt{\tau}Q(X,Y_i)d\tau \\
&  = \frac{n}{\sqrt{\bar{\tau}}}-2\sqrt{\bar{\tau}}R
-\int_0^{\bar{\tau}}\sqrt{\tau}\(\frac{\tau}{\bar{\tau}}\)Q(X)d\tau \nn\\
& =\frac{n}{\sqrt{\bar{\tau}}}-2\sqrt{\bar{\tau}}R
-\frac{1}{\bar{\tau}}K. \nn
\end{align} %\eqno(3.2.10)
Moreover, by Proposition 3.2.3, the equality in (3.2.10) holds
everywhere if and only if for each $(q,\bar{\tau})$ and any
shortest $\mathcal{L}$-geodesic $\gamma$ on $[0,\bar{\tau}]$
connecting $p$ and $q$, and for any unit vector $Y$ at
$\tau=\bar{\tau}$, the extended vector field $Y(\tau)$ along
$\gamma$ by the ODE (3.2.8) must be an $\mathcal{L}$-Jacobian
field. When $Y_i(\tau),$ $i=1,\ldots,n$ are $\mathcal{L}$-Jacobian
fields along $\gamma$, we have
\begin{align*}
& \frac{d}{d\tau}\langle Y_i(\tau),Y_j(\tau)\rangle \\
&  = 2\Ric(Y_i,Y_j)+\langle\nabla_XY_i,Y_j\rangle
+\langle Y_i,\nabla_XY_j\rangle\\
& =2\Ric(Y_i,Y_j)+\left\langle\nabla_{Y_i}
\(\frac{1}{2\sqrt{\bar{\tau}}}\nabla L\),Y_j\right\rangle
+\left\langle Y_i,\nabla_{Y_j}
\(\frac{1}{2\sqrt{\bar{\tau}}}\nabla L\)\right\rangle\\
&  = 2\Ric(Y_i,Y_j)+\frac{1}{\sqrt{\bar{\tau}}}{\rm
Hess}_L(Y_i,Y_j)
\end{align*}
and then by (3.2.9),
$$
2\Ric(Y_i,Y_j)+\frac{1}{\sqrt{\bar{\tau}}}{\rm Hess}_L(Y_i,Y_j)
=\frac{1}{\bar{\tau}}\delta_{ij},\ \ \ \text{at}\ \
\tau=\bar{\tau}.
$$
Therefore the equality in (3.2.10) holds everywhere if and only if
we are on a gradient shrinking soliton with
$$
R_{ij}+\frac{1}{2\sqrt{\bar\tau}}\nabla_i\nabla_jL
=\frac{1}{2\bar\tau}g_{ij}.
$$
\end{proof}
%$$\eqno\#$$ \vskip 0.3cm

In summary, from (3.2.6), (3.2.7) and Corollary 3.2.4, we have
$$   \left\{
       \begin{array}{lll}
  \frac{\partial L}{\partial \bar{\tau}}
=2\sqrt{\bar{\tau}}R-\frac{L}{2\bar{\tau}}
  +\frac{K}{\bar{\tau}},\\[4mm]
  |\nabla L|^2=-4\bar{\tau}R+\frac{2}{\sqrt{\bar{\tau}}}L
-\frac{4}{\sqrt{\bar{\tau}}}K,\\[4mm]
  \Delta
  L\leq-2\sqrt{\bar{\tau}}R+\frac{n}{\sqrt{\bar{\tau}}}
-\frac{K}{\bar{\tau}},
       \end{array}
    \right.
$$
in the sense of distributions.

Now the \textbf{Li-Yau-Perelman distance} \index{Li-Yau-Perelman
distance} $l=l(q,\bar{\tau})$ is defined by
$$l
(q,\bar{\tau})=L(q,\bar{\tau})/2\sqrt{\bar{\tau}}.
$$
We thus have the following

%\vskip 0.2cm\noindent{\bf Lemma 3.2.5}\ \emph{
% CZ-BASELINE source-line:5795 source-kind:lemma source-id:3.2.5
\begin{lemma}[identities for Li–Yau–Perelman distance]
  \label{lem:cz-identities-for-liyauperelman-distance}
  \dcref{Lemma 3.2.5}
  \group{Cao--Zhu Comparison Theorems for Perelman's Reduced Volume}
  \source{cao-zhu}
For the Li-Yau-Perelman distance $l(q,\bar{\tau})$ defined above,
we have
\begin{align}
\frac{\partial l}{\partial  \bar{\tau}}
&=-\frac{l}{\bar{\tau}}+R+\frac{1}{2\bar{\tau}^{3/2}}K,\\
 %\eqno (3.2.11)
|\nabla l|^2&=-R+\frac{l}{\bar{\tau}}-\frac{1}{\bar{\tau}^{3/2}}K,\\
%\eqno(3.2.12)
\Delta l
&\leq-R+\frac{n}{2\bar{\tau}}-\frac{1}{2\bar{\tau}^{3/2}}K,
 % \eqno(3.2.13)
\end{align}
in the sense of distributions. Moreover, the equality in
$(3.2.13)$ holds if and only if we are on a gradient shrinking
soliton.
%$$\eqno \#$$ \vskip 0.3cm
\end{lemma}

As the first consequence, we derive the following upper bound on
the minimum of $l(\cdot,\tau)$ for every $\tau$ which will be
useful in proving the no local collapsing theorem in the next
section.

%\vskip 0.2cm \noindent {\bf Corollary 3.2.6} \ \emph{
% CZ-BASELINE source-line:5820 source-kind:corollary source-id:3.2.6
\begin{corollary}[conjugate heat inequality for reduced density]
  \label{cor:cz-conjugate-heat-inequality-for-reduced-density}
  \dcref{Corollary 3.2.6}
  \group{Cao--Zhu Comparison Theorems for Perelman's Reduced Volume}
  \source{cao-zhu}
Let $g_{ij}(\tau)$, $\tau\geq 0$, be a family of metrics evolving
by the Ricci flow $\frac{\partial}{\partial \tau}g_{ij}=2R_{ij}$
on a compact $n$-dimensional manifold $M$. Fix a point $p$ in $M$
and let $l(q,\tau)$ be the Li-Yau-Perelman distance from $(p,0)$.
Then for all $\tau$,
$$
\min\{l(q,\tau)\ |\ \ q\in M\}\leq\frac{n}{2}.
$$
\end{corollary}

%\vskip 0.1cm \noindent {\bf Proof.} \ \
\begin{proof}
Let
$$
\bar{L}(q,\tau)=4\tau l(q,\tau).
$$
Then, it follows from (3.2.11) and (3.2.13) that
$$
\frac{\partial \bar{L}}{\partial \tau}=4\tau
R+\frac{2K}{\sqrt{\tau}},
$$
and
$$
\Delta \bar{L}\leq-4\tau R+2n-\frac{2K}{\sqrt{\tau}}.
$$
Hence
$$
\frac{\partial \bar{L}}{\partial \tau}+\Delta \bar{L}\leq 2n.
$$
Thus, by a standard maximum principle argument,
$\min\{\bar{L}(q,\tau)-2n\tau| \ q\in M\}$ is nonincreasing and
therefore $\min\{\bar{L}(q,\tau)|\ \ q\in M\}\leq 2n\tau$.
%$$\eqno\#$$\vskip 0.3cm
\end{proof}

As another consequence of Lemma 3.2.5, we obtain
$$
\frac{\partial l}{\partial \bar{\tau}}-\Delta l+|\nabla l|^2-R
+\frac{n}{2\bar{\tau}}\geq 0.
$$
or equivalently
$$
\(\frac{\partial}{\partial \bar{\tau}}-\Delta +R\)
\((4\pi\bar{\tau})^{-\frac{n}{2}}\exp(-l)\)\leq 0.
$$
If $M$ is compact, we define {\bf Perelman's reduced
volume}\index{Perelman's reduced volume} by
$$
\tilde{V}(\tau)=\int_M(4\pi\tau)^{-\frac{n}{2}}
\exp(-l(q,\tau))dV_{\tau}(q),
$$
where $dV_{\tau}$ denotes the volume element with respect to the
metric $g_{ij}(\tau)$. Note that Perelman's reduced volume
resembles the expression in Huisken's monotonicity formula for the
mean curvature flow \cite{Hu90}. It follows, from the above
computation, that
\begin{align*}
& \frac{d}{d\bar{\tau}}\int_M(4\pi\bar{\tau})^{-\frac{n}{2}}
\exp(-l(q,\bar{\tau}))dV_{\bar{\tau}}(q)\\
&  = \int_M\left[\frac{\partial}{\partial \bar{\tau}}
((4\pi\bar{\tau})^{-\frac{n}{2}}\exp(-l(q,\bar{\tau})))
+R(4\pi\bar{\tau})^{-\frac{n}{2}}\exp(-l(q,\bar{\tau}))\right]
dV_{\bar{\tau}}(q)\\
&  \leq \int_M\Delta((4\pi\bar{\tau})^{-\frac{n}{2}}
\exp(-l(q,\bar{\tau})))dV_{\bar{\tau}}(q)\\
&  = 0.
\end{align*}
This says that if $M$ is compact, then Perelman's reduced volume
$\tilde{V}(\tau)$ is nonincreasing in $\tau$; moreover, the
monotonicity is strict unless we are on a gradient shrinking
soliton.

\medskip
In order to define and to obtain the monotonicity of  Perelman's
reduced volume for a complete noncompact manifold, we need to
formulate the monotonicity of Perelman's reduced volume in a local
version. This local version is very important and will play a
crucial role in the analysis of the Ricci flow with surgery in
Chapter 7.

\medskip
We define the {\bf $\mathcal{L}$-exponential map (with parameter
$\bar \tau$)}\index{$\mathcal{L}$-exponential map!with parameter
$\bar \tau$} ${\mathcal{L}}\exp(\bar \tau): T_pM \rightarrow M $
as follows: for any $X\in T_pM$, we set
$$
{\mathcal{L}}\exp_X(\bar{\tau})=\gamma(\bar{\tau})
$$
where $\gamma(\tau)$ is the ${\mathcal{L}}$-geodesic, starting at
$p$ and having $X$ as the limit of $\sqrt{\tau}\dot{\gamma}(\tau)$
as $\tau\rightarrow 0^+$. The associated Jacobian of the
$\mathcal{L}$-exponential map is called {\bf
$\mathcal{L}$-Jacobian}\index{$\mathcal{L}$-Jacobian}. We denote
by $\mathcal{J}(\tau)$ the $\mathcal{L}$-Jacobian  of
${\mathcal{L}}\exp(\tau):$ $T_pM\rightarrow M $. We can now deduce
an estimate for the $\mathcal{L}$-Jacobian as follows.

Let $q={\mathcal{L}}\exp_X(\bar{\tau})$ and $\gamma(\tau)$,
$\tau\in[0,\bar{\tau}]$, be the shortest ${\mathcal{L}}$-geodesic
connecting $p$ and $q$ with
$\sqrt{\tau}\dot{\gamma}(\tau)\rightarrow X$ as $\tau\rightarrow
0^+$. For any vector $v\in T_pM$, we consider the family of
${\mathcal{L}}$-geodesics:
$$
\gamma_s(\tau)={\mathcal{L}}\exp_{(X+sv)}(\tau),  \ \ \ \
0\leq\tau\leq\bar{\tau},\ \  s\in(-\epsilon,\epsilon).
$$
The associated variation vector field $V(\tau)$,
$0\leq\tau\leq\bar{\tau}$, is an $\mathcal{L}$-Jacobian field with
$V(0)=0$ and $V(\tau)=({\mathcal{L}}\exp_X(\tau))_*(v)$.

Let $v_1,\dots,v_n$ be $n$ linearly independent vectors in $T_pM$.
Then
$$
V_i(\tau)=({\mathcal{L}}\exp_X(\tau))_*(v_i),  \ \ \ \ i=1,2,
\ldots,n,
$$
are $n$ $\mathcal{L}$-Jacobian fields along $\gamma(\tau)$,
$\tau\in[0,\bar{\tau}]$. The $\mathcal{L}$-Jacobian
$\mathcal{J}(\tau)$ is given by
$$
{\mathcal{J}(\tau)}=|V_1(\tau)\wedge\dots\wedge
V_n(\tau)|_{g_{ij}(\tau)}/|v_1\wedge\dots\wedge v_n|.
$$

Now for fixed $b\in(0,\bar{\tau})$, we can choose linearly
independent vectors $v_1,\ldots,v_n\in T_pM$ such that $\langle
V_i(b),V_j(b)\rangle_{g_{ij}(b)}=\delta_{ij}$. We compute
\begin{align*}
&\frac{d}{d\tau}\mathcal{J}^2 \\
&  =\frac{2}{|v_1\wedge\cdots\wedge
v_n|^2}\sum\limits_{j=1}^n\langle V_1\wedge\cdots\wedge
\nabla_XV_j\wedge\cdots\wedge V_n,V_1\wedge\cdots\wedge
V_n\rangle_{g_{ij}(\tau)}\\
&\; +\frac{2}{|v_1\wedge\cdots\wedge
v_n|^2}\sum\limits_{j=1}^n\langle V_1\wedge\cdots\wedge
\Ric(V_j,\cdot)\wedge\cdots\wedge V_n,V_1\wedge\cdots\wedge
V_n\rangle_{g_{ij}(\tau)}.
\end{align*}
At $\tau=b$,
$$
\frac{d}{d\tau}\mathcal{J}^2(b)=\frac{2}{|v_1\wedge\cdots\wedge
v_n|^2}\sum\limits_{j=1}^n(\langle\nabla_XV_j,V_j\rangle_{g_{ij}(b)}
+\Ric(V_j,V_j)).
$$
Thus,
\begin{align*}
\frac{d}{d\tau}\log\mathcal{J}(b) &
=\sum\limits_{j=1}^n(\langle\nabla_XV_j,V_j\rangle_{g_{ij}(b)}
+\Ric(V_j,V_j))\\
& =\sum\limits_{j=1}^n\(\left\langle\nabla_{V_j}
\(\frac{1}{2\sqrt{b}}\nabla L\),V_j\right\rangle_{g_{ij}(b)}
+\Ric(V_j,V_j)\)\\
&  = \frac{1}{2\sqrt{b}}\(\sum\limits_{j=1}^n{\rm Hess}_L(V_j,V_j)\)+R\\
&  = \frac{1}{2\sqrt{b}}\Delta L+R.
\end{align*}
Therefore, in view of Corollary 3.2.4, we obtain the following
estimate for $\mathcal{L}$-Jacobian: \be
\frac{d}{d\tau}\log{\mathcal{J}}(\tau)
\leq\frac{n}{2\tau}-\frac{1}{2\tau^{3/2}}K\
\ \ \text{on} \ \ \ [0,\bar{\tau}].  %\eqno(3.2.14)
\ee On the other hand, by the definition of the Li-Yau-Perelman
distance and (3.2.4), we have
\begin{align}
\frac{d}{d\tau} l(\tau)
&  =-\frac{1}{2\tau}l+\frac{1}{2\sqrt{\tau}}\frac{d}{d\tau}L \\
&  =-\frac{1}{2\tau}l+\frac{1}{2\sqrt{\tau}}(\sqrt{\tau}(R+|X|^2)) \nn\\
&  =-\frac{1}{2\tau^{3/2}}K.\nn
\end{align}   %\eqno(3.2.15)
Here and in the following we denote by $l(\tau) =
l(\gamma(\tau),\tau)$.  Now the combination of (3.2.14) and
(3.2.15) implies the following important {\bf Jacobian comparison
theorem}\index{Jacobian comparison theorem} of Perelman \cite{P1}.

%\vskip 0.2cm\noindent {\bf Theorem 3.2.7} (Perelman's Jacobian
%comparison theorem) \ \emph{
% CZ-BASELINE source-line:5997 source-kind:theorem source-id:3.2.7
\begin{theorem}[Perelman’s reduced Jacobian comparison]
  \label{thm:cz-perelmans-reduced-jacobian-comparison}
  \dcref{Theorem 3.2.7}
  \group{Cao--Zhu Comparison Theorems for Perelman's Reduced Volume}
  \source{cao-zhu}
Let $g_{ij}(\tau)$ be a family of complete solutions to the Ricci
flow $\frac{\partial}{\partial \tau}g_{ij}=2R_{ij}$ on a manifold
$M$ with bounded curvature. Let $\gamma:[0,\bar\tau]\rightarrow M$
be a shortest $\mathcal{L}$-geodesic starting from a fixed point
$p$. Then {\bf Perelman's reduced volume element}
$$
(4\pi\tau)^{-\frac{n}{2}}\exp(-l(\tau)){\mathcal{J}(\tau)}
$$
is nonincreasing in $\tau$ along $\gamma$.\index{Perelman's
reduced volume! element}
\end{theorem}
%$$\eqno \#$$ \vskip 0.3cm

We now show how to integrate Perelman's reduced volume element
over $T_pM$ to deduce the following monotonicity result of
Perelman \cite{P1}.

%\vskip 0.2cm \noindent {\bf Theorem 3.2.8}
%(Monotonicity of Perelman's reduced volume)
% CZ-BASELINE source-line:6017 source-kind:theorem source-id:3.2.8
\begin{theorem}[monotonicity of reduced volume]
  \label{thm:cz-monotonicity-of-reduced-volume}
  \dcref{Theorem 3.2.8}
  \group{Cao--Zhu Comparison Theorems for Perelman's Reduced Volume}
  \source{cao-zhu}
Let $g_{ij}$ be a family of complete metrics evolving by the Ricci
flow $\frac{\partial}{\partial \tau}g_{ij}=2R_{ij}$ on a manifold
$M$ with bounded curvature. Fix a point $p$ in $M$ and let
$l(q,\tau)$ be the reduced distance from $(p,0)$. Then
\begin{itemize}
\item[(i)] Perelman's reduced volume
$$
\tilde{V}(\tau)
=\int_M(4\pi\tau)^{-\frac{n}{2}}\exp(-l(q,\tau))dV_{\tau}(q)
$$
is finite and nonincreasing in $\tau$; \item[(ii)] the
monotonicity is strict unless we are on a gradient shrinking
soliton.
\end{itemize}
\end{theorem}

%\vskip 0.1cm\noindent {\bf Proof.}
\begin{proof}
  \uses{thm:cz-global-curvature-derivative-estimates,thm:cz-perelmans-reduced-jacobian-comparison}
For any $v\in T_pM$ we can find an $\mathcal{L}$-geodesic
$\gamma(\tau)$, starting at $p$, with
$\lim\limits_{\tau\rightarrow0^+}\sqrt{\tau}\dot{\gamma}(\tau)=v$.
Recall that $\gamma(\tau)$ satisfies the $\mathcal{L}$-geodesic
equation
$$
\nabla_{\dot{\gamma}(\tau)}\dot{\gamma}(\tau)-\frac{1}{2}\nabla
R+\frac{1}{2\tau}\dot{\gamma}(\tau)+2\Ric(\dot{\gamma}(\tau),\cdot)=0.
$$
Multiplying this equation by $\sqrt{\tau}$, we get \be
\frac{d}{d\tau}(\sqrt{\tau}\dot{\gamma})-\frac{1}{2}\sqrt{\tau}\nabla
R+2\Ric(\sqrt{\tau}\dot{\gamma}(\tau),\cdot)=0.%\eqno(3.2.16)
\ee Since the curvature of the metric $g_{ij}(\tau)$ is bounded,
it follows from Shi's derivative estimate (Theorem 1.4.1) that
$|\nabla R|$ is also bounded for small $\tau>0$. Thus by
integrating (3.2.16), we have \be
|\sqrt{\tau}\dot{\gamma}(\tau)-v|\leq C\tau(|v|+1) %\eqno(3.2.17)
\ee for $\tau$ small enough and for some positive constant $C$
depending only the curvature bound.

Let $v_1,\ldots,v_n$ be $n$ linearly independent vectors in $T_pM$
and let
$$
V_i(\tau)=(\mathcal{L}\exp_v(\tau))_*(v_i)
=\frac{d}{ds}|_{s=0}\mathcal{L}\exp_{(v+sv_i)}(\tau),\ \
i=1,\ldots,n.
$$
The $\mathcal{L}$-Jacobian $\mathcal{J}(\tau)$ is given by
$$
\mathcal{J}(\tau)=|V_1(\tau)\wedge\cdots\wedge
V_n(\tau)|_{g_{ij}(\tau)}/|v_1\wedge\cdots\wedge v_n|
$$
By (3.2.17), we see that
$$
\left|\sqrt{\tau}\frac{d}{d\tau}\mathcal{L}
\exp_{(v+sv_i)}(\tau)-(v+sv_i)\right| \leq C\tau (|v|+|v_i|+1)
$$
for $\tau$ small enough and all $s\in(-\epsilon,\epsilon)$ (for
some $\epsilon>0$ small) and $i=1,\ldots,n$. This implies that
$$
\lim\limits_{\tau\rightarrow0^+}\sqrt{\tau}\dot{V}_i(\tau) =v_i,\
\ \ i=1,\ldots,n,
$$
so we deduce that \be
\lim\limits_{\tau\rightarrow0^+}\tau^{-\frac{n}{2}}\mathcal{J}(\tau)=1.
%\eqno(3.2.18)
\ee Meanwhile, by using (3.2.17), we have
\begin{displaymath}
\begin{split}
l(\tau)&  =\frac{1}{2\sqrt{\tau}}\int_0^\tau\sqrt{\tau}
(R+|\dot{\gamma}(\tau)|^2)d\tau\\
&  \rightarrow |v|^2\ \ \text{as}\ \tau\rightarrow0^+.
\end{split}
\end{displaymath}
Thus \be
l(0)=|v|^2.%\eqno(3.2.19)
\ee Combining (3.2.18) and (3.2.19) with Theorem 3.2.7, we get
\begin{displaymath}
\begin{split}
\tilde{V}(\tau)
&  =\int_M(4\pi\tau)^{-\frac{n}{2}}\exp(-l(q,\tau))dV_{\tau}(q) \\
&  \leq\int_{T_pM}(4\pi\tau)^{-\frac{n}{2}}
\exp(-l(\tau))\mathcal{J}(\tau)|_{\tau=0}dv\\
&  =(4\pi)^{-\frac{n}{2}}\int_{\mathbb{R}^n}\exp(-|v|^2)dv\\
&  <+\infty.
\end{split}
\end{displaymath}
This proves that Perelman's reduced volume is always finite and
hence well defined. Now the monotonicity assertion in (i) follows
directly from Theorem 3.2.7.

For the assertion (ii), we note that the equality in (3.2.13)
holds everywhere when the monotonicity of  Perelman's reduced
volume is not strict. Therefore we have completed the proof of the
theorem.
\end{proof}
%$$\eqno\#$$


\section{No Local Collapsing Theorem I}

In this section we apply the monotonicity of Perelman's reduced
volume in Theorem 3.2.8 to prove Perelman's  {\bf no local
collapsing theorem I}\index{no local collapsing theorem I}(cf.
section 7.3 and section 4 of \cite{P1}), which is extremely
important not only because it gives a local injectivity radius
estimate in terms of local curvature bound but also it will
survive the surgeries in Chapter 7.


%\vskip 0.2cm \noindent {\bf Definition 3.3.1}.
% CZ-BASELINE source-line:6127 source-kind:definition source-id:3.3.1
\begin{definition}[kappa-noncollapsing at a scale]
  \label{def:cz-kappa-noncollapsing-at-a-scale}
  \dcref{Definition 3.3.1}
  \group{Cao--Zhu No Local Collapsing Theorem I}
  \source{cao-zhu}
Let $\kappa$, $r$ be two positive constants and let $g_{ij}(t), 0
\leq t <T,$ be a solution to the Ricci flow on an $n$-dimensional
manifold $M$. We call the solution $g_{ij}(t)$
\textbf{$\kappa$-noncollapsed} at $(x_0,t_0)\in M \times [0,T)$ on
the scale $r$ \index{$\kappa$-noncollapsed} if it satisfies the
following property: whenever
$$
|Rm|(x,t)\leq r^{-2}
$$
for all $x \in B_{t_0}(x_0,r)$ and $t \in [t_0-r^2,t_0]$, we have
$$Vol_{t_0}(B_{t_0}(x_0,r)) \geq \kappa r^n.$$ Here $B_{t_0}(x_0,r)$ is
the geodesic ball centered at $x_0\in M $ and of radius $r$ with
respect to the metric $g_{ij}(t_0)$.
\end{definition}

Now we are ready to state the {\bf no local collapsing theorem
I}\index{no local collapsing theorem I} of Perelman \cite{P1}.

%\vskip 0.2cm \noindent {\bf Theorem 3.3.2} (No local collapsing
%theorem I)\
% CZ-BASELINE source-line:6148 source-kind:theorem source-id:3.3.2
\begin{theorem}[no local collapsing on a finite time interval]
  \label{thm:cz-no-local-collapsing-on-a-finite-time-interval}
  \dcref{Theorem 3.3.2}
  \group{Cao--Zhu No Local Collapsing Theorem I}
  \source{cao-zhu}
Given any metric $g_{ij}$ on an $n$-dimensional compact manifold
$M$. Let $g_{ij}(t)$ be the solution to the Ricci flow on $[0,T)$,
with $T<+\infty$, starting at $g_{ij}$. Then there exist positive
constants $\kappa$ and $\rho_0$ such that for any $t_0\in[0,T)$
and any point $x_0\in M$, the solution $g_{ij}(t)$ is
$\kappa$-noncollapsed at $(x_0,t_0)$ on all scales less than
$\rho_0$.
\end{theorem}

%\vskip 0.1cm \noindent {\bf Proof.}
\begin{proof}
  \uses{thm:cz-local-curvature-derivative-estimates,thm:cz-perelmans-reduced-jacobian-comparison,cor:cz-conjugate-heat-inequality-for-reduced-density}
We argue by contradiction. Suppose that there are sequences
$p_k\in M$, $t_k\in[0,T)$ and $r_k\rightarrow 0$ such that \be
|Rm|(x,t)\leq r_k^{-2} %\eqno (3.3.1)
\ee for $x\in B_k=B_{t_k}(p_k,r_k)$ and $t_k-r_k^2\leq t\leq t_k$,
but \be \epsilon_k=r_k^{-1}Vol_{t_k}(B_k)^{\frac{1}{n}}\rightarrow
0 \
\ \text{as} \ k\rightarrow \infty. %\eqno(3.3.2)
\ee

Without loss of generality, we may assume that $t_k\rightarrow T$
as $k\rightarrow +\infty$. \vskip 0.1cm \noindent Let
$\bar{\tau}(t)=t_k-t$, $p=p_k$ and
$$
\tilde{V_k}(\bar{\tau})=\int_M(4\pi\bar{\tau})^{-\frac{n}{2}}
\exp(-l(q,\bar{\tau}))dV_{t_k-\bar{\tau}}(q),
$$
where $l(q,\bar\tau)$ is the Li-Yau-Perelman distance with respect
to $p=p_k$.

\medskip
{\it Step} 1. \ We first want to show that for $k$ large enough,
$$
\tilde{V_k}(\epsilon_kr_k^2)\leq2\epsilon_k^\frac{n}{2}.
$$

For any $v\in T_pM$ we can find an $\mathcal{L}$-geodesic
$\gamma(\tau)$ starting at $p$ with $\lim\limits_{\tau\rightarrow
0}\sqrt{\tau}\dot{\gamma}(\tau)$ $=v$. Recall that $\gamma(\tau)$
satisfies the equation (3.2.16). It follows from assumption
(3.3.1) and Shi's local derivative estimate (Theorem 1.4.2) that
$|\nabla R|$ has a bound in the order of ${1}/{r_k^3}$ for $t\in
[t_k-\epsilon_kr_k^2,t_k]$. Thus by integrating (3.2.16) we see
that for $\tau\leq\epsilon_kr_k^2$ satisfying the property that
$\gamma(\sigma)\in B_k$ as long as $\sigma<\tau$, there holds \be
|\sqrt{\tau}\dot{\gamma}(\tau)-v|\leq C\epsilon_k(|v|+1)
%\eqno(3.3.3)
\ee where $C$ is some positive constant depending only on the
dimension. Here we have implicitly used the fact that the metric
$g_{ij}(t)$ is equivalent for $x\in B_k$ and $t\in
[t_k-\epsilon_kr_k^2,t_k]$. In fact since $\frac{\partial
g_{ij}}{\partial t}=-2R_{ij}$ and $|Rm|\leq r_k^{-2}$ on
$B_k\times [t_k-r_k^2,t_k]$, we have \be
e^{-2\epsilon_k}g_{ij}(x,t_k)\leq g_{ij}(x,t)\leq
e^{2\epsilon_k}g_{ij}(x,t_k), %\eqno (3.3.4)
\ee for $x\in B_k$ and $t\in [t_k-\epsilon_kr_k^2,t_k]$.\vskip
0.3cm Suppose $v\in T_pM$ with
$|v|\leq\frac{1}{4}\epsilon_k^{-\frac{1}{2}}$. Let $\tau\leq
\epsilon_kr_k^2$ such that $\gamma(\sigma)\in B_k$ as long as
$\sigma<\tau$, where $\gamma$ is the ${\mathcal{L}}$-geodesic
starting at $p$ with $\lim\limits_{\tau\rightarrow
0}\sqrt{\tau}\dot{\gamma}(\tau)=v$. Then, by (3.3.3) and (3.3.4),
for $k$ large enough,
\begin{align*}
d_{t_k}(p_k,\gamma(\tau))
&  \leq \int_0^\tau|\dot{\gamma}(\sigma)|_{g_{ij}(t_k)}d\sigma\\
&  < \frac{1}{2}\epsilon_k^{-\frac{1}{2}}
\int_0^\tau\frac{d\sigma}{\sqrt{\sigma}}\\
&  = \epsilon_k^{-\frac{1}{2}}\sqrt{\tau}\\
&  \leq r_k.
\end{align*}
This shows that for $k$ large enough, \be
{\mathcal{L}}\exp_{\{|v|\leq\frac{1}{4}\epsilon_k^{-1/2}\}}
(\epsilon_kr_k^2)\subset
B_k=B_{t_k}(p_k,r_k). %\eqno(3.3.5)
\ee We now estimate the integral of $\tilde{V_k}(\epsilon_kr_k^2)$
as follows,
\begin{align}
\tilde{V_k}(\epsilon_kr_k^2) &
=\int_M(4\pi\epsilon_kr_k^2)^{-\frac{n}{2}}
\exp(-l(q,\epsilon_kr_k^2))dV_{t_k-\epsilon_kr_k^2}(q)\\
& =\int\limits_{{\mathcal{L}}\exp_{\{|v|\leq\frac{1}{4}
\epsilon_k^{-1/2}\}}(\epsilon_kr_k^2)}
(4\pi\epsilon_kr_k^2)^{-\frac{n}{2}}
\exp(-l(q,\epsilon_kr_k^2))dV_{t_k-\epsilon_kr_k^2}(q)\nn\\
&\quad +\!\int\limits_{M\setminus{\mathcal{L}}
\exp_{\{|v|\leq\frac{1}{4}\epsilon_k^{-1/2}\}}(\epsilon_kr_k^2)}\!\!
(4\pi\epsilon_kr_k^2)^{-\frac{n}{2}}
\exp(-l(q,\epsilon_kr_k^2))dV_{t_k-\epsilon_kr_k^2}(q).\nn
\end{align}   %\eqno(3.3.6)
We observe that for each $q\in B_k$,
$$
L(q,\epsilon_kr^2_k)
=\int_0^{\epsilon_kr_k^2}\sqrt{\tau}(R+|\dot{\gamma}|^2)d\tau
\geq-C(n)r_k^{-2}(\epsilon_kr_k^2)^{\frac{3}{2}}
=-C(n)\epsilon_k^{\frac{3}{2}}r_k,
$$
hence $l(q,\epsilon_kr^2_k)\geq -C(n)\epsilon_k.$ Thus, the first
term on the RHS of (3.3.6) can be estimated by
\begin{align}
&\int\limits_{\mathcal{L}\exp_{\{|v|\leq\frac{1}{4}
\epsilon_k^{-1/2}\}}(\epsilon_kr^2_k)}
(4\pi\epsilon_kr^2_k)^{-\frac{n}{2}}
\exp(-l(q,\epsilon_kr_k^2))dV_{t_k-\epsilon_kr^2_k}(q)\\
&\leq e^{n\epsilon_k}\int\limits_{B_k}
(4\pi\epsilon_kr^2_k)^{-\frac{n}{2}}
\exp(-l(q,\epsilon_kr_k^2))dV_{t_k}(q)\nn\\
&\leq e^{n\epsilon_k}(4\pi)^{-\frac{n}{2}}\cdot
e^{C(n)\epsilon_k}\cdot
\epsilon_k^{-\frac{n}{2}}\cdot(r_k^{-n}\Vol_{t_k}(B_k))\nn\\
&=e^{(n+C(n))\epsilon_k}(4\pi)^{-\frac{n}{2}}
\cdot\epsilon_k^\frac{n}{2},\nn
\end{align}  %\eqno(3.3.7)
where we have also used (3.3.5) and (3.3.4).

Meanwhile, by using (3.2.18), (3.2.19) and the Jacobian Comparison
Theorem 3.2.7, the second term on the RHS of (3.3.6) can be
estimated as follows
\begin{align}
&\int\limits_{M\setminus\mathcal{L}\exp_{\{|v|\leq\frac{1}{4}
\epsilon_k^{-\frac{1}{2}}\}}(\epsilon_kr^2_k)}
(4\pi\epsilon_kr^2_k)^{-\frac{n}{2}}
\exp(-l(q,\epsilon_kr_k^2))dV_{t_k-\epsilon_kr^2_k}(q)\\
&\leq \int\limits_{\{|v|>\frac{1}{4}\epsilon_k^{-\frac{1}{2}}\}}
(4\pi\tau)^{-\frac{n}{2}}
\exp(-l(\tau))\mathcal{J}(\tau)|_{\tau=0}dv \nn\\
&=(4\pi)^{-\frac{n}{2}}\int\limits_{\{|v|>\frac{1}{4}
\epsilon_k^{-\frac{1}{2}}\}}\exp(-|v|^2)dv \nn\\
&\leq \epsilon_k^{\frac{n}{2}},\nn
\end{align} %\eqno(3.3.8)
for $k$ sufficiently large. Combining (3.3.6)-(3.3.8), we finish
the proof of Step 1.

\medskip
{\it Step} 2. \ We next want to show
$$
\tilde{V_k}(t_k)=(4\pi
t_k)^{-\frac{n}{2}}\int_M\exp(-l(q,t_k))dV_0(q)>C'
$$
for all $k$, where $C'$ is some positive constant independent of
$k$.

It suffices to show the Li-Yau-Perelman distance $l(\cdot, t_k)$
is uniformly bounded from above on $M$. By Corollary 3.2.6 we know
that the minimum of $l(\cdot,\tau)$ does not exceed $\frac{n}{2}$
for each $\tau>0$. Choose $q_k\in M$ such that the minimum of
$l(\cdot,t_k-\frac{T}{2})$ is attained at $q_k$. We now construct
a path $\gamma: [0,t_k]\rightarrow M$ connecting $p_k$ to any
given point $q\in M$ as follows: the first half path
$\gamma|_{[0,t_k-\frac{T}{2}]}$ connects $p_k$ to $q_k$ so that
$$
l\(q_k,t_k-\frac{T}{2}\) =\frac{1}{2\sqrt{t_k-\frac{T}{2}}}
\int_0^{t_k-\frac{T}{2}}\sqrt{\tau}
(R+|\dot{\gamma}(\tau)|^2)d\tau\leq\frac{n}{2}
$$
and the second half path $\gamma|_{[t_k-\frac{T}{2},t_k]}$ is a
shortest geodesic connecting $q_k$ to $q$ with respect to the
initial metric $g_{ij}(0)$. Then, for any $q\in M^n$,
\begin{align*}
l(q,t_k)&  = \frac{1}{2\sqrt{t_k}}L(q,t_k)\\
& \leq \frac{1}{2\sqrt{t_k}}\(\int_0^{t_k-\frac{T}{2}}
+\int_{t_k-\frac{T}{2}}^{t_k}\)
\sqrt{\tau}(R+|\dot{\gamma}(\tau)|^2)d\tau\\
& \leq \frac{1}{2\sqrt{t_k}}\(n\sqrt{t_k-\frac{T}{2}}
+\int_{t_k-\frac{T}{2}}^{t_k}
\sqrt{\tau}(R+|\dot{\gamma}(\tau)|^2)d\tau\)\\
&  \leq  C
\end{align*}
for some constant $C>0$, since all geometric quantities in
$g_{ij}$ are uniformly bounded when $t\in[0,\frac{T}{2}]$ (or
equivalently, $\tau\in[t_k-\frac{T}{2},t_k]$).

Combining Step 1 with Step 2, and using the monotonicity of
$\tilde{V_k}(\tau)$, we get
$$
C'<\tilde{V_k}(t_k)\leq\tilde{V_k}(\epsilon_kr_k^2) \leq
2\epsilon_k^\frac{n}{2}\to 0
$$
as $k\to \infty$. This gives the desired contradiction. Therefore
we have proved the theorem.
\end{proof}
%$$\eqno\#$$ \vskip 0.3cm

The above no local collapsing theorem I says that if $|Rm|\leq
r^{-2}$ on the parabolic ball $\{(x,t)\ |\ d_{t_0}(x,x_0)\leq r, \
t_0-r^2\leq t\leq t_0\}$, then the volume of the geodesic ball
$B_{t_0}(x_0,r)$ (with respect to the metric $g_{ij}(t_0)$) is
bounded from below by $\kappa r^n$. In \cite{P1}, Perelman used
the monotonicity of the $\mathcal{W}$-functional (defined by
(1.5.9)) to obtain a stronger version of the no local collapsing
theorem, where the curvature bound assumption on the parabolic
ball is replaced by that on the geodesic ball $B_{t_0}(x_0,r)$.
The following result, called {\bf no local collapsing theorem
I$'$}\index{no local collapsing theorem I$'$}, gives a further
extension where the bound on the curvature tensor is replaced by
the bound on the scalar curvature only.

%\vskip 0.2cm\noindent {\bf Theorem 3.3.3} (No local collapsing
%theorem I$'$)
% CZ-BASELINE source-line:6348 source-kind:theorem source-id:3.3.3
\begin{theorem}[scale-local no collapsing]
  \label{thm:cz-scale-local-no-collapsing}
  \dcref{Theorem 3.3.3}
  \group{Cao--Zhu No Local Collapsing Theorem I}
  \source{cao-zhu}
Suppose $M$ is a compact Riemannian manifold, and $g_{ij}(t)$,
$0\leq t<T<+\infty$, is a solution to the Ricci flow. Then there
exists a positive constant $\kappa$ depending only the initial
metric and $T$ such that for any $(x_0,t_0)\in M\times(0,T)$ if
$$
R(x,t_0)\leq r^{-2},\ \ \ \forall x\in B_{t_0}(x_0,r)
$$
with $0<r\leq\sqrt{T}$, then we have
$$
{\rm Vol}_{t_0}(B_{t_0}(x_0,r))\geq\kappa r^n.
$$
\end{theorem}

%\vskip 0.1cm \noindent {\bf Proof.}
\begin{proof}
  \uses{cor:cz-mu-monotonicity-and-shrinking-breathers}
We will prove the assertion
$$
%(\ast)_a\ \ \ \ \ \
{\rm Vol}_{t_0}(B_{t_0}(x_0,a))\geq\kappa a^n\leqno{(\ast)_a}
$$
for all $0<a\leq r$. Recall that
$$
\mu(g_{ij},\tau)=\inf\left\{\mathcal{W}(g_{ij},f,\tau)\ \Big|\
\int_M(4\pi\tau)^{-\frac{n}{2}}e^{-f}dV=1\right\}.
$$
Set
$$
\mu_0=\inf\limits_{0\leq\tau\leq2T}\mu(g_{ij}(0),\tau)>-\infty.
$$
By Corollary 1.5.9, we have
\begin{align}
\mu(g_{ij}(t_0),b)&  \geq\mu(g_{ij}(0),t_0+b)\\
                  &  \geq\mu_0 \nn
\end{align}     %\eqno (3.3.9)
for $0<b\leq r^2$. Let $0<\zeta\leq1$ be a positive smooth
function on $\mathbb{R}$ where $\zeta(s)=1$ for
$|s|\leq\frac{1}{2}$, $|\zeta'|^2/\zeta\leq20$ everywhere, and
$\zeta(s)$ is very close to zero for $|s|\geq1$. Define a function
$f$ on $M$ by
$$
(4\pi r^2)^{-\frac{n}{2}}e^{-f(x)}=e^{-c}(4\pi
r^2)^{-\frac{n}{2}}\zeta\(\frac{d_{t_0}(x,x_0)}{r}\),
$$
where the constant $c$ is chosen so that $\int_M(4\pi
r^2)^{-\frac{n}{2}}e^{-f}dV_{t_0}=1$. Then it follows from (3.3.9)
that
\begin{align}
\mathcal{W}(g_{ij}(t_0),f,r^2) &  =\int_M[r^2(|\nabla f|^2
+R)+f-n]
(4\pi r^2)^{-\frac{n}{2}}e^{-f}dV_{t_0}\\
                 &  \geq\mu_0.\nn
\end{align} %\eqno (3.3.10)
Note that
\begin{displaymath}
\begin{split}
   1&  =\int_M(4\pi
   r^2)^{-\frac{n}{2}}e^{-c}\zeta\(\frac{d_{t_0}(x,x_0)}{r}\)dV_{t_0}\\
    &  \geq\int_{B_{t_0}(x_0,\frac{r}{2})}(4\pi
    r^2)^{-\frac{n}{2}}e^{-c}dV_{t_0}\\
    &  =(4\pi
   r^2)^{-\frac{n}{2}}e^{-c}\Vol_{t_0}\(B_{t_0}\(x_0,\frac{r}{2}\)\).
\end{split}
\end{displaymath}
By combining with (3.3.10) and the scalar curvature bound, we have
\begin{displaymath}
\begin{split}
   c&  \geq-\int_M\(\frac{(\zeta')^2}{\zeta}
-\log\zeta\cdot\zeta\)e^{-c}(4\pi r^2)^{-\frac{n}{2}}
dV_{t_0}+(n-1)+\mu_0\\
    &  \geq-2(20+e^{-1})e^{-c}(4\pi
    r^2)^{-\frac{n}{2}}\Vol_{t_0}(B_{t_0}(x_0,r))+(n-1)+\mu_0\\
    &  \geq-2(20+e^{-1})
      \frac{{\rm Vol}_{t_0}(B_{t_0}(x_0,r))}{\Vol_{t_0}(B_{t_0}
(x_0,\frac{r}{2}))}+(n-1)+\mu_0,
\end{split}
\end{displaymath}
where we used the fact that $\zeta(s)$ is very close to zero for
$|s|\geq1$. Note also that
$$
2\int_{B_{t_0}(x_0,r)}e^{-c}(4\pi r^2)^{-\frac{n}{2}}dV_{t_0}
\geq\int_M(4\pi r^2)^{-\frac{n}{2}}e^{-f}dV_{t_0}=1.
$$
Let us set
$$\kappa=\min\left\{\frac{1}{2}\exp(-2(20+e^{-1})3^{-n}+(n-1)+\mu_0),\
\frac{1}{2}\alpha_n\right\}
$$
where $\alpha_n$ is the volume of the unit ball in $\mathbb{R}^n$.
Then we obtain
\begin{displaymath}
\begin{split}
  {\rm Vol}_{t_0}(B_{t_0}(x_0,r))
   &  \geq\frac{1}{2}e^c(4\pi r^2)^{\frac{n}{2}}\\
   &  \geq\frac{1}{2}(4\pi)^{\frac{n}{2}}
\exp(-2(20+e^{-1})3^{-n}+(n-1)+\mu_0)\cdot  r^n\\
   &  \geq\kappa r^n
\end{split}
\end{displaymath}
provided Vol$_{t_0}(B_{t_0}(x_0,\frac{r}{2}))
\geq3^{-n}Vol_{t_0}(B_{t_0}(x_0,r))$.

Note that the above argument also works for any smaller radius
$a\leq r$. Thus we have proved the following assertion: \be
{\rm Vol}_{t_0}(B_{t_0}(x_0,a))\geq\kappa a^n  %\eqno(3.3.11)
\ee whenever $a\in(0,r]$ and ${\rm
Vol}_{t_0}(B_{t_0}(x_0,\frac{a}{2})) \geq3^{-n}{\rm
Vol}_{t_0}(B_{t_0}(x_0,a))$.

Now we argue by contradiction to prove the assertion $(\ast)_a$
for any $a\in(0,r]$. Suppose $(\ast)_a$ fails for some
$a\in(0,r]$. Then by (3.3.11) we have
\begin{displaymath}
\begin{split}
   \Vol_{t_0}(B_{t_0}(x_0,\frac{a}{2}))
&  <3^{-n}\Vol_{t_0}(B_{t_0}(x_0,a))\\
   &  <3^{-n}\kappa a^n\\
   &  <\kappa\(\frac{a}{2}\)^n.
\end{split}
\end{displaymath}
This says that $(\ast)_{\frac{a}{2}}$ would also fail. By
induction, we deduce that
$$
\Vol_{t_0}\(B_{t_0}\(x_0,\frac{a}{2^k}\)\)<\kappa\(\frac{a}{2^k}\)^n\
\ \ \text{for all } k\geq1.
$$
This is a contradiction since $\lim\limits_{k\rightarrow\infty}
\Vol_{t_0}\(B_{t_0}\(x_0,\frac{a}{2^k}\)\)/\(\frac{a}{2^k}\)^n=\alpha_n$.
%$$\eqno\#$$ \vskip 1cm
\end{proof}


\section{No Local Collapsing Theorem II}

By inspecting the arguments in the previous section, one can see
that if the injectivity radius of the initial metric is uniformly
bounded from below, then the no local collapsing theorem I of
Perelman also holds for complete solutions with bounded curvature
on a complete noncompact manifold. In this section we will use a
cut-off argument to extend the no local collapsing theorem to any
complete solution with bounded curvature. In some sense, the
second no local collapsing theorem of Perelman \cite{P1} gives a
good relative estimate of the volume element for the Ricci flow.

We first need the following useful lemma which contains two
assertions. The first one is a parabolic version of the Laplacian
comparison theorem (where the curvature sign restriction in the
ordinary Laplacian comparison is essentially removed in the Ricci
flow). The second one is a generalization of a result of Hamilton
(Theorem 17.2 in \cite{Ha95F}), where it was derived by an
integral version of Bonnet-Myers' theorem.

%\vskip 0.2cm \noindent {\bf Lemma 3.4.1} (Perelman \cite{P1})\ \
% CZ-BASELINE source-line:6500 source-kind:lemma source-id:3.4.1
\begin{lemma}[distance distortion under Ricci flow]
  \label{lem:cz-distance-distortion-under-ricci-flow}
  \dcref{Lemma 3.4.1}
  \group{Cao--Zhu No Local Collapsing Theorem II}
  \source{cao-zhu}
Let $g_{ij}(x,t)$ be a solution to the Ricci flow on an
$n$-dimensional manifold $M$ and denote by $d_t(x,x_0)$ the
distance between $x$ and $x_0$ with respect to the metric
$g_{ij}(t)$.
\begin{itemize}
\item[(i)] Suppose $\Ric(\cdot,t_0)\leq(n-1)K$ on
$B_{t_0}(x_0,r_0)$ for some $x_0\in M$ and some positive constants
K and $r_0$. Then the distance function $d(x,t)=d_t(x,x_0)$
satisfies, at $t=t_0$ and outside $B_{t_0}(x_0,r_0)$, the
differential inequality:
$$
\frac{\partial}{\partial t}d-\Delta d
\geq-(n-1)\(\frac{2}{3}Kr_0+r_0^{-1}\).
$$
\item[(ii)] Suppose $\Ric(\cdot,t_0)\leq (n-1)K$ on
$B_{t_0}(x_0,r_0)\bigcup B_{t_0}(x_1,r_0)$ for some $x_0,x_1\in M$
and some positive constants K and $r_0$.  Then, at $t=t_0$,
$$
\frac{d}{dt}d_t(x_0,x_1)\geq-2(n-1)\(\frac{2}{3}Kr_0+r_0^{-1}\).
$$
\end{itemize}
\end{lemma}

%\vskip 0.1cm \noindent {\bf Proof.}
\begin{proof}
  \uses{prop:cz-l-jacobi-comparison-inequality}
Let $\gamma: [0,d(x,t_0)]\rightarrow M$ be a shortest normal
geodesic from $x_0$ to $x$ with respect to the metric
$g_{ij}(t_0)$. As usual, we may assume that $x$ and $x_0$ are not
conjugate to each other in the metric $g_{ij}(t_0)$, otherwise we
can understand the differential inequality in the barrier sense.
Let $X=\dot{\gamma}(0)$ and let $\{X,e_1,\ldots,e_{n-1}\}$ be an
orthonormal basis of $T_{x_0}M$. Extend this basis parallel along
$\gamma$ to form a parallel orthonormal basis
$\{X(s),e_1(s),\ldots,e_{n-1}(s)\}$ along $\gamma$.

\medskip
(i) Let $X_i(s),\ i=1,\ldots,n-1$, be the Jacobian fields along
$\gamma$ such that $X_i(0)=0$ and $X_i(d(x,t_0))=e_i(d(x,t_0))$
for $i=1,\ldots,n-1$. Then it is well-known that (see for example
\cite{ScY})
$$
\Delta d_{t_0}(x,x_0)
=\sum_{i=1}^{n-1}\int_0^{d(x,t_0)}(|\dot{X}_i|^2-R(X,X_i,X,X_i))ds
$$
(in Proposition 3.2.3 we actually did this for the more
complicated $\mathcal{L}$-distance function).

Define vector fields $Y_i,\ i=1,\ldots,n-1$, along $\gamma$ as
follows:
$$
Y_i(s)=\begin{cases}
\frac{s}{r_0}e_i(s),&   \text{if }\;s\in[0,r_0],\\
e_i(s),&   \text{if }\;s\in[r_0,d(x,t_0)].\end{cases}
$$
which have the same value as the corresponding Jacobian fields
$X_i(s)$ at the two end points of $\gamma$. Then by using the
standard index comparison theorem (see for example \cite{CE}) we
have
\begin{align*}
\Delta d_{t_0}(x,x_0) &  =
\sum\limits_{i=1}^{n-1}\int_0^{d(x,t_0)}
(|\dot{X}_i|^2-R(X,X_i,X,X_i))ds\\
&  \leq \sum\limits_{i=1}^{n-1}\int_0^{d(x,t_0)}(|\dot{Y}_i|^2
-R(X,Y_i,X,Y_i))ds\\
&  = \int_0^{r_0}\frac{1}{r_0^2}(n-1-s^2\Ric(X,X))ds
+\int_{r_0}^{d(x,t_0)}(-\Ric(X,X))ds\\
&  = -\int_\gamma \Ric(X,X)+\int^{r_0}_0\(\frac{(n-1)}{r_0^2}
+\(1-\frac{s^2}{r_0^2}\)\Ric(X,X)\)ds\\
&  \leq -\int_\gamma \Ric(X,X)+(n-1)\(\frac{2}{3}Kr_0+r_0^{-1}\).
\end{align*}
On the other hand,
\begin{align*}
\frac{\partial}{\partial t}d_t(x,x_0) & = \frac{\partial}{\partial
t}
\int_0^{d(x,t_0)}\sqrt{g_{ij}X^iX^j}ds\\
&  = -\int_\gamma \Ric(X,X)ds.
\end{align*}
Hence we obtain the desired differential inequality. \vskip 3mm

\medskip
(ii) The proof is divided into three cases.

\medskip
{\it Case} (1): $d_{t_0}(x_0,x_1)\geq 2r_0$.

Define vector fields $Y_i,\ i=1,\ldots,n-1$, along $\gamma$ as
follows:
$$
Y_i(s)=\begin{cases}
\frac{s}{r_0}e_i(s),&   \mbox{ if }\ s\in[0,r_0],\\[4mm]
e_i(s),&   \mbox{ if }\ s\in[r_0,d(x_1,t_0)],\\[4mm]
\frac{d(x_1,t_0)-s}{r_0}e_i(s),&   \mbox{ if }\
s\in[d(x_1,t_0)-r_0,d(x_1,t_0)].\end{cases}
$$
Then by the second variation formula, we have
$$
\sum\limits_{i=1}^{n-1}\int_0^{d(x_1,t_0)}R(X,Y_i,X,Y_i)ds
\leq\sum\limits_{i=1}^{n-1}\int_0^{d(x_1,t_0)}|\dot{Y}_i|^2ds,
$$
which implies
\begin{align*}
& \int_0^{r_0}\frac{s^{2}}{r_0^2}\Ric(X,X)ds
+\int_{r_0}^{d(x,t_0)-r_0}\Ric(X,X)ds\\
&\quad+\int_{d(x_1,t_0)-r_0}^{d(x_1,t_0)}
\(\frac{d(x_1,t_0)-s}{r_0}\)^{2}\Ric(X,X)ds \leq
\frac{2(n-1)}{r_0}.
\end{align*}
Thus
\begin{align*}
& \frac{d}{dt}(d_t(x_0,x_1))\\
& \geq-\int_0^{r_0}\(1-\frac{s^{2}}{r_0^2}\)\Ric(X,X)ds \\
&\quad-\int_{d(x_1,t_0)-r_0}^{d(x_1,t_0)}
\(1-\(\frac{d(x_1,t_0)-s}{r_0}\)^{2}\)\Ric(X,X)ds
%&\quad
- \frac{2(n-1)}{r_0}\\
& \geq-2(n-1)\(\frac{2}{3}Kr_0+r_0^{-1}\).
\end{align*}

\smallskip
{\it Case} (2): $\frac{2}{\sqrt{\frac{2K}{3}}}\leq
d_{t_0}(x_0,x_1)\leq2r_0$.

In this case, letting $r_1=\frac{1}{\sqrt{\frac{2K}{3}}}$ and
applying case (1) with $r_0$ replaced by $r_1,$ we get
\begin{align*}
\frac{d}{dt}(d_t(x_0,x_1))
&  \geq -2(n-1)\(\frac{2}{3}Kr_1+r_1^{-1}\)\\
&  \geq -2(n-1)\(\frac{2}{3}Kr_0+r_0^{-1}\).
\end{align*}

\medskip
{\it Case} (3): $ d_{t_0}(x_0,x_1)\leq
\min\Big\{\frac{2}{\sqrt{\frac{2K}{3}}}, 2r_0\Big\}$.

In this case,
$$
\int_0^{d(x_1,t_0)}\Ric(X,X)ds \leq (n-1)K
\frac{2}{\sqrt{\frac{2K}{3}}}=(n-1)\sqrt{6K},
$$
and
$$
2(n-1)\(\frac{2}{3}Kr_0+r_0^{-1}\)\geq (n-1)\sqrt{\frac{32}{3}K}.
$$
This proves the lemma.
%$$\eqno \#$$ \vskip 0.3cm
\end{proof}

The following result, called the {\bf no local collapsing theorem
II}\index{no local collapsing theorem II}, was obtained by
Perelman in \cite{P1}.

%%\vskip 0.2cm \noindent 3.4.2

% CZ-BASELINE source-line:6654 source-kind:theorem source-id:3.4.2
\begin{theorem}[no local collapsing for complete flows]
  \label{thm:cz-no-local-collapsing-for-complete-flows}
  \dcref{Theorem 3.4.2}
  \group{Cao--Zhu No Local Collapsing Theorem II}
  \source{cao-zhu}
For any $A >0$ there exists $\kappa=\kappa(A)>0$ with the
following property: if $g_{ij}(t)$ is a complete solution to the
Ricci flow on $0\leq t\leq r_0^2$ with bounded curvature and
satifying
$$
|Rm|(x,t)\leq r_0^{-2}\ \ \ on \ B_0(x_0,r_0)\times [0,r_0^2]
$$
and
$$
\Vol_0(B_0(x_0,r_0))\geq A^{-1}r_0^n,
$$
then $g_{ij}(t)$ is $\kappa$-noncollapsed on all scales less than
$r_0$ at every point $(x,r_0^2)$ with $d_{r_0^2}(x,x_0)\leq Ar_0$.
\end{theorem}

%\noindent {\bf Proof.}
\begin{proof}
  \uses{thm:cz-no-local-collapsing-on-a-finite-time-interval,lem:cz-distance-distortion-under-ricci-flow}
{}From the evolution equation of the Ricci flow, we know that the
metrics $g_{ij}(\cdot,t)$ are equivalent to each other on
$B_0(x_0,r_0)\times [0,r_0^2]$. Thus, without loss of generality,
we may assume that the curvature of the solution is uniformly
bounded for all $t \in [0,r_0^2]$ and all points in
$B_t(x_0,r_0)$. Fix a point $(x,r_0^2) \in M\times \{r_0^2\}$. By
scaling we may assume $r_0=1$. We may also assume $d_1(x,x_0)=A$.
Let $p=x$, $\bar{\tau}=1-t$, and consider Perelman's reduced
volume
$$
\tilde{V}(\bar{\tau}) =\int_M(4\pi\bar{\tau})^{-\frac{n}{2}}
\exp(-l(q,\bar{\tau}))dV_{1-\bar{\tau}}(q),
$$
where
\begin{multline*}
l(q,\bar{\tau})=\inf\bigg\{\frac{1}{2\sqrt{\bar{\tau}}}
\int_0^{\bar{\tau}}\sqrt{\tau}(R+|\dot{\gamma}|^2)d\tau
\mid\ \gamma: [0,\bar{\tau}]\rightarrow M\\
\mbox{with}\ \gamma(0)=p,\ \gamma(\bar{\tau})=q \bigg\}
\end{multline*}
is the Li-Yau-Perelman distance. We argue by contradiction.
Suppose for some $0<r<1$ we have
$$
|Rm|(y,t)\leq r^{-2}
$$
whenever $y\in B_{1}(x,r)$ and $1-r^2\leq t\leq 1$, but
$\epsilon=r^{-1}\Vol_{1}(B_1(x,r))^{\frac{1}{n}}$ is very small.
Then arguing as in the proof of the no local collapsing theorem I
(Theorem 3.3.2), we see that Perelman's reduced volume
$$
\tilde{V}(\epsilon r^2)\leq 2\epsilon^\frac{n}{2}
$$
On the other hand, from the monotonicity of Perelman's reduced
volume we have
$$
(4\pi)^{-\frac{n}{2}}\int_M\exp(-l(q,1))dV_0(q)
=\tilde{V}(1)\leq\tilde{V}(\epsilon r^2).
$$
Thus once we bound the function $l(q,1)$ over $B_0(x_0,1)$ from
above, we will get the desired contradiction and will prove the
theorem.

For any $q\in B_0(x_0,1)$, exactly as in the proof of the no local
collapsing theorem I, we choose a path $\gamma: [0,1]\rightarrow
M$ with $\gamma(0)=x,\ \gamma(1)=q$, $\gamma(\frac{1}{2})=y\in
B_\frac{1}{2}(x_0,\frac{1}{10})$ and $\gamma (\tau) \in
B_{1-\tau}(x_0,1)$ for $\tau \in [\frac{1}{2},1]$ such that
$$
{\mathcal{L}}(\gamma|_{[0,\frac{1}{2}]})
=2\sqrt{\frac{1}{2}}l\(y,\frac{1}{2}\)\ \ \(=L\(y,\frac{1}{2}\)\).
$$
Now ${\mathcal{L}}(\gamma|_{[\frac{1}{2},1]})
=\int_\frac{1}{2}^1\sqrt{\tau}(R(\gamma(\tau),1-\tau)+
|\dot{\gamma}(\tau)|^2_{g_{ij}(1-\tau)})d\tau$ is bounded from
above by a uniform constant since all geometric quantities in
$g_{ij}$ are uniformly bounded on $\{(y,t)\ |\ t\in[0,{1}/{2}],
y\in B_t(x_0,1)\}$ (where $t\in [0,1/2]$ is equivalent to
$\tau\in[{1}/{2},1]$).  Thus all we need is to estimate the
minimum of $l(\cdot,\frac{1}{2})$, or equivalently
$\bar{L}(\cdot,\frac{1}{2})=4\frac{1}{2}l(\cdot,\frac{1}{2})$, in
the ball $B_\frac{1}{2}(x_0,\frac{1}{10})$.

Recall that $\bar{L}$ satisfies the differential inequality \be
\frac{\partial \bar{L}}{\partial \tau}+\Delta\bar{L}\leq 2n.
%\eqno(3.4.1)
\ee We will use this in a maximum principle argument. Let us
define
$$
h(y,t)=\phi(d(y,t)-A(2t-1))\cdot(\bar{L}(y,1-t)+2n+1)
$$
where $d(y,t)=d_t(y,x_0)$, and $\phi$ is a function of one
variable, equal to 1 on $(-\infty,\frac{1}{20})$, and rapidly
increasing to infinity on $(\frac{1}{20},\frac{1}{10})$ in such a
way that: \be
2\frac{(\phi')^2}{\phi}-\phi''\geq(2A+100n)\phi'-C(A)\phi
%\eqno(3.4.2)
\ee for some constant $C(A)<+\infty$. The existence of such a
function $\phi$ can be justified as follows: put
$v=\frac{\phi'}{\phi}$, then the condition (3.4.2) for $\phi$ can
be written as
$$
3v^2-v'\geq(2A+100n)v-C(A)
$$
which can be solved for $v$.

Since the scalar curvature $R$ evolves by
$$
\frac{\partial R}{\partial t}=\Delta R+2|Rc|^2\geq\Delta
R+\frac{2}{n}R^2,
$$
we can apply the maximum principle as in Chapter 2 to deduce
$$
R(x,t)\geq-\frac{n}{2t}\ \ \text{for}\ t\in(0,1]\ \text{and} \
x\in M.
$$
Thus for $\bar{\tau}=1-t\in[0,\frac{1}{2}]$,
\begin{align*}
\bar{L}(\cdot,\bar{\tau}) &
=2\sqrt{\bar{\tau}}\int_0^{\bar{\tau}}
\sqrt{\tau}(R+|\dot{\gamma}|^2)d\tau\\
&  \geq 2\sqrt{\bar{\tau}}\int_0^{\bar{\tau}}
\sqrt{\tau}\(-\frac{n}{2(1-\tau)}\)d\tau\\
&  \geq 2\sqrt{\bar{\tau}}\int_0^{\bar{\tau}}
\sqrt{\tau}(-n)d\tau\\
&  > -2n.
\end{align*}
That is \be \bar{L}(\cdot,1-t)+2n+1\geq1,\ \ \text{for}\
t\in\left[\frac{1}{2},1\right]. %\eqno(3.4.3)
\ee Clearly $\min\limits_{y\in M}h(y,\frac{1}{2})$ is achieved by
some $y\in B_{\frac{1}{2}}(x_0,\frac{1}{10})$ and \be
\min\limits_{y\in M}h(y,1)\leq h(x,1)=2n+1.  %\eqno(3.4.4)
\ee We compute
\begin{align*}
\(\frac{\partial}{\partial t}-\Delta\)h
& =\(\frac{\partial}{\partial t}-\Delta\)\phi\cdot(\bar{L}(y,1-t)+2n+1)\\
&\quad+\phi\cdot\(\frac{\partial}{\partial
t}-\Delta\)\bar{L}(y,1-t)
 -2\langle\nabla\phi,\nabla\bar{L}(y,1-t)\rangle\\
&  =\(\phi'\left[\(\frac{\partial}{\partial t}-\Delta\)d-2A\right]
-\phi''|\nabla d|^2\)\cdot(\bar{L}+2n+1)\\
&\quad +\phi\cdot\(-\frac{\partial}{\partial \tau}-\Delta\)
\bar{L}-2\langle\nabla\phi,\nabla\bar{L}\rangle\\
&  \geq \(\phi'\left[\(\frac{\partial}{\partial
t}-\Delta\)d-2A\right]-\phi''\)\cdot(\bar{L}+2n+1)\\
&\quad-2n\phi-2\langle\nabla\phi,\nabla\bar{L}\rangle
\end{align*}
by using (3.4.1). At a minimizing point of $h$ we have
$$
\frac{\nabla\phi}{\phi}=-\frac{\nabla\bar{L}}{(\bar{L}+2n+1)}.
$$
Hence
$$
-2\langle\nabla\phi,\nabla\bar{L}\rangle
=2\frac{|\nabla\phi|^2}{\phi}(\bar{L}+2n+1)
=2\frac{(\phi')^2}{\phi}(\bar{L}+2n+1).$$ Then at the minimizing
point of $h$, we compute
\begin{align*}
\(\frac{\partial}{\partial t}-\Delta\)h& \geq
\(\phi'\left[\(\frac{\partial}{\partial
t}-\Delta\)d-2A\right]-\phi''\)
\cdot(\bar{L}+2n+1) \\
&\quad-2n\phi+2\frac{(\phi')^2}{\phi}(\bar{L}+2n+1)\\
&  \geq \(\phi'\left[\(\frac{\partial}{\partial t}
-\Delta\)d-2A\right]-\phi''\)
\cdot(\bar{L}+2n+1) \\
&\quad-2nh+2\frac{(\phi')^2}{\phi}(\bar{L}+2n+1)
\end{align*}
for $t\in[\frac{1}{2},1]$ and
$$
\Delta h\geq0.
$$
Let us denote by $h_{\min}(t)=\min\limits_{y\in M}h(y,t)$. By
applying Lemma 3.4.1(i) to the set where $\phi'\neq0$, we further
obtain
\begin{align*}
\frac{d}{dt}h_{\min} &  \geq
(\bar{L}+2n+1)\cdot\left[\phi'(-100n-2A)
-\phi''+2\frac{(\phi')^2}{\phi}\right]-2nh_{\min}\\
&  \geq -(2n+C(A))h_{\min},\quad \text{for }\;
t\in[\frac{1}{2},1].
\end{align*}
This implies that $h_{\min}(t)$ cannot decrease too fast. By
combining (3.4.3) and (3.4.4) we get the required estimate for the
minimum $\bar{L}(\cdot,\frac{1}{2})$ in the ball
$B_{\frac{1}{2}}(x_0,\frac{1}{10})$.

Therefore we have completed the proof of the theorem.
%$$\eqno \#$$ \vskip 1cm
\end{proof}

%\bigskip
